\chapter{Higher-order miniKanren: the literature}\label{ch:lkanren}

\providecommand{\lkclo}[2]{\langle #1;\,#2\rangle}
\providecommand{\lkq}{\stackrel{?}{=}}
\providecommand{\lkJust}{\mathsf{Just}}
\providecommand{\lkNothing}{\mathsf{Nothing}}
\providecommand{\lkyes}{{\ttfamily ✓}}
\providecommand{\lkno}{{\ttfamily ✗}}

The previous chapters described what λProlog is. It has a logic, higher-order hereditary Harrop
formulas (\cref{ch:hh}). It has a term language, simply typed λ-terms used as λ-tree syntax
(\cref{ch:lambda,ch:lambdatree}). It has a unifier, Huet's procedure or its decidable core, pattern
unification (\cref{ch:hounif}). And it has a style of programming (\cref{ch:apps}). This chapter
changes the point of view. Several people have tried to put these ideas into a miniKanren, and
several systems implement λProlog efficiently. What did they actually do, and what went wrong?

The user named one paper as the key reference: Ma and Friedman, \emph{A New Higher-order
Unification Algorithm for λKanren} \cite{MaFriedman2021}. We present it in full in
\cref{sec:lkanren-mf21}, followed by a critical assessment against pattern unification and nominal
unification (\cref{sec:lkanren-assess}). Then we survey the rest of the literature. The earlier
λKanren paper \cite{MaLuFriedman2020} comes first (\cref{sec:lkanren-mk20}). Then nominal unification
and αKanren \cite{UrbanPittsGabbay2004,ByrdFriedman2007} (\cref{sec:lkanren-nominal}). Then
eigenvariables, negation and ∀ inside miniKanren
\cite{HemannByrdFriedman2013,Moiseenko2019,JinEtAl2021} (\cref{sec:lkanren-eigen}). Last come the
techniques used by real λProlog systems, Teyjus, ELPI and Makam (\cref{sec:lkanren-impl}). The chapter
ends with a comparison table of the design options (\cref{sec:lkanren-compare}). It does
\emph{not} choose between them: the choice is made in \cref{ch:design}.

\section{The map of the territory}\label{sec:lkanren-map}

The book gives the logic and an operational semantics. It has the search rules \rulename{and},
\rulename{augment}, \rulename{generic} and \rulename{backchain}, plus unification of λ-terms. A
miniKanren that wants to host this must answer six concrete questions. The rest of the chapter is
organized around them.

\begin{enumerate}[label=\textbf{Q\arabic*.}]
  \item \textbf{Representation of λ-terms.} The options are named binders, de Bruijn indices, de
    Bruijn levels, locally nameless terms, nominal atoms, or explicit substitutions. The goal is to
    make α-equivalence cheap and to avoid β as far as possible.
  \item \textbf{Which unification runs inside \scm{=°}.} There are five candidates: first-order (today),
    nominal (αKanren), Miller's pattern unification $\Llam$ (λKanren 2020, Teyjus~2, ELPI), Huet's full
    pre-unification, or the \enquote{static closures} of Ma and Friedman 2021.
  \item \textbf{How to implement \rulename{generic} ($\forall x.G$).} You need fresh eigen-constants
    plus a \emph{scope check}. The check can use an explicit quantifier prefix
    \cite{Miller1991,MaLuFriedman2020}, numeric universe indices or levels
    \cite{NadathurJayaramanKwon1995,DunchevEtAl2015}, or an absento-like constraint over the
    variables in scope \cite{HemannByrdFriedman2013}.
  \item \textbf{How to implement \rulename{augment} ($D\imp G$).} A hypothetical clause must be visible
    to every relation call made during the proof of $G$, and only there. In a shallow embedding this
    clashes with the lexical scope of host-language procedures.
  \item \textbf{What to do outside the decidable fragment.} You can fail, delay the equation as a
    constraint, or enumerate unifiers.
  \item \textbf{How the new features meet the other constraints} \scm{≠°}, \scm{absent°} and the
    type tags. Eigen-constants and disequality do not mix naively.
\end{enumerate}

The historical arc, in the notes' order: Miller 1991 ($\Llam$, mixed prefixes, raising and pruning)
\cite{Miller1991}; Nipkow 1993 (functional, de Bruijn) \cite{Nipkow1993}; Nadathur, Jayaraman and Kwon
1995 (universe indices) \cite{NadathurJayaramanKwon1995}; Dowek, Hardin, Kirchner and Pfenning 1996
(explicit substitutions) \cite{DowekHardinKirchnerPfenning1996}; Urban, Pitts and Gabbay 2004 (nominal
unification) \cite{UrbanPittsGabbay2004}; Nadathur and Linnell 2005 (on-the-fly raising)
\cite{NadathurLinnell2005}; αKanren 2007 \cite{ByrdFriedman2007}; ELPI 2015 \cite{DunchevEtAl2015};
then the miniKanren papers of 2019--2021
\cite{Moiseenko2019,MaLuFriedman2020,MaFriedman2021,JinEtAl2021}.

% =====================================================================================================
\section{Ma and Friedman 2021: static closures}\label{sec:lkanren-mf21}

\subsection{What kind of paper this is}

The paper is a ten-page workshop paper, presented at the miniKanren Workshop 2021. Its abstract states
the thesis (p.\,1):
\begin{quote}
\enquote{λKanren, adapted from λProlog, extends miniKanren with higher-order logic programming.
The existing higher-order unification algorithms, however, are neither sound nor complete w.r.t.\ an
“intuitive” understanding of α-equivalence. This paper shows these gaps and proposes a new method to
handle α-equivalence in higher-order unification.}
\end{quote}
It opens provocatively: \enquote{Is there a world where true is equal to false? Indeed, there is, if
higher-order unification and miniKanren meet each other} (p.\,1).

The background section recalls the one notion from \cref{ch:hounif} that the whole paper argues
about, in its own words (p.\,1): \enquote{For an application form $Fa$, if $F$ is a unification
variable, then $a$ must be a list of zero or more distinct bound variables. Such restricted application
forms are called \emph{patterns} and we hereafter refer to the problem identified by Miller as
\emph{pattern unification}.} The related work (p.\,7--8) adds Huet's vocabulary: a term is \emph{rigid}
\enquote{if its head symbol is a bound variable or a constant}, and \emph{flex} \enquote{if its head symbol
is a unification variable}. Keep in mind one fact for later: in every query of §2--4, the unknown
$X$ sits under a λ but is \emph{not} applied to anything. It is a pattern with zero arguments, so
pattern unification lets it depend on none of the bound variables.

\begin{pitfall}[a preview, not an algorithm]
The conclusion says: \enquote{This paper shows a preview of static closures} (p.\,9). The paper has
\textbf{no algorithm listing, no inference rules for unification, no theorems and no proofs}. Its
formal content is two grammars, three operations on scopes and three worked examples. Its code is a
set of REPL transcripts. Soundness and completeness are explicitly future work: \enquote{We
anticipate extending the formulation of higher-order unification with static closures and derive a
sound and complete algorithm for α-equivalence} (p.\,9). The rules in \cref{sec:lkanren-recon} below
are our own \emph{reconstruction}, not the paper's, and we flag them as such. (Where this chapter
says \enquote{the notes}, it means the unpublished working notes from which this survey was written.)
\end{pitfall}

The paper has seven sections. §1 gives background. §2 asks \enquote{grafting or substitution?}
(p.\,2--4). §3 shows how one can \enquote{turn false to true} (p.\,4). §4 introduces static closures
(p.\,5--6). §5 describes a shallow-embedding implementation of the language (p.\,6--7). §6 is related
work (p.\,7--8) and §7 the conclusion. The paper uses two box styles for examples: \enquote{solid
boxes for examples that are complete and dash boxes for examples that are incomplete or somewhat
questionable} (footnote 2, p.\,2). In every \scm{(run 1 X ...)} of the paper, \scm{X} is simply
existentially quantified at the top, outside every λ (footnote 1, p.\,2).

\subsection{Grafting versus substitution}\label{sec:lkanren-graft}

First-order miniKanren already has a word for what unification produces. For

\begin{scheme}
> (run 1 X
    (== X 'cat))
'(cat)
\end{scheme}

\noindent the paper says \enquote{internally, the answer is a grafting \dots\ written as $\{X\mapsto
\texttt{'cat}\}$} (p.\,2).

\begin{definition}[Grafting and substitution {\cite[p.\,2--3]{MaFriedman2021}}]\label{def:lkanren-graft}
A \emph{grafting} $\{X\mapsto t\}$ is a purely textual replacement of the unification variable $X$
by $t$. It is \emph{not} capture-avoiding: grafting $X\mapsto a$ into $\abs{a}{X}$ yields
$\abs{a}{a}$, so the name $a$ gets captured. A \emph{substitution} is the capture-avoiding
instantiation of standard higher-order unification: \enquote{Standard higher-order unification relies
on substitution, but not grafting, to instantiate unification variables. Substitution is
capture-avoiding, whereas grafting is not} (p.\,3).
\end{definition}

In first-order unification the two notions coincide, because there are no binders. They part ways
as soon as terms have binders. The paper follows a \enquote{novice} through three queries. The first
one, α-equivalence of closed terms, is unproblematic:

\begin{scheme}
> (run 1 X
    (== (λ (a) a) (λ (b) b)))
'(_0)
\end{scheme}

\noindent (\texttt{\_0} means that $X$ is unrestricted). Then, \enquote{out of curiosity}, the novice
puts the query variable under a binder and asks \enquote{Does this query succeed?}:

\begin{scheme}
> (run 1 X
    (== (λ (a) X) (λ (b) b)))
'()
\end{scheme}

\noindent \enquote{Surprisingly enough, with higher-order unification, there is no answer.} The novice
expects $X = a$. The reason is the \emph{scope restriction}.

\begin{definition}[Scope restriction {\cite[p.\,3]{MaFriedman2021}}]\label{def:lkanren-scope}
\enquote{Unification variables are ∃-quantified and lambda variables are ∀-quantified. A scope is a
list of those quantifiers. Further, a ∀-quantified variable may instantiate (or occurs in the
instantiation of) an ∃-quantified variable, only if the ∀ variable precedes the ∃ variable in the
scope.} As a side condition on a binding in a mixed prefix $Q$:
\[
  \inferrule{X \text{ is } \exists\text{-quantified in } Q \\ \forall y\in FV_\forall(t).\; y \prec_Q X}
            {X\mapsto t \text{ is allowed}}
\]
\end{definition}

For \lp{λa.X =? λb.b} the scope is \scm{'((∃ X) (∀ a))}. The variable $a$ comes \emph{after} $X$, so
it cannot instantiate $X$. This is exactly λProlog's behaviour: it is the eigenvariable condition of
\rulename{generic} (\cref{ch:hh}), and in unification it is what raising and pruning enforce
(\cref{ch:hounif}). The paper draws the opposite conclusion (p.\,3): \enquote{It is fair to say that
the existing pattern unification algorithms are not complete w.r.t.\ the natural sense of
α-equivalence.} It also says (p.\,3): \enquote{For beginners who are not familiar with higher-order
unification, α-equivalence is only decided by the binder positions. But pattern unification imposes
an appendage limitation by taking unification variables into account.}

A conjunction shows why the restriction has a reason to exist (dashed box, p.\,3):

\begin{scheme}
> (run 1 X
    (== (λ (a) X) (λ (b) b))
    (== (λ (c) X) (λ (d) d)))
\end{scheme}

\noindent \enquote{Associating X with a cannot work, neither does c. If we text-replace X with either,
then the other one fails.} The paper's first proposal is to stop treating λ-variables as names:
\enquote{We first propose to treat unification variables and lambda variables differently. In
particular, lambda variables are effectively constants if they are represented by de Bruijn levels}
(p.\,3).

\begin{definition}[de Bruijn levels and indices {\cite[footnote 3, p.\,4]{MaFriedman2021}}]\label{def:lkanren-levels}
\enquote{There are two ways to implement de Bruijn's nameless syntax, indices and levels. Both use
numbers to represent variables} \cite{deBruijn1972}. A \emph{level} counts binders from the
\emph{outside}: the outermost binder is $0$. An \emph{index} counts binders from the \emph{inside}:
the nearest binder is $0$.
\end{definition}

\begin{example}
In $\abs{c}{\abs{d}{c}}$, the variable $c$ has level $0$ and index $1$, and $d$ has level $1$ and
index $0$. The term is $\lambda.\lambda.\,0$ with levels and $\lambda.\lambda.\,1$ with indices. The
paper writes the level $\ell$ as \texttt{d}$\ell$. So \texttt{d0} names the outermost enclosing
binder, whatever it is called, and the answer \texttt{'(d0)} reads: \enquote{these two terms can be made
α-equal, if X corresponds to de Bruijn level 0} (p.\,3).
\end{example}

With levels, the two problematic queries succeed (solid boxes, p.\,3):

\begin{scheme}
> (run 1 X
    (== (λ (a) X) (λ (b) b)))
'(d0)
> (run 1 X
    (== (λ (a) X) (λ (b) b))
    (== (λ (c) X) (λ (d) d)))
'(d0)
\end{scheme}

\noindent \enquote{If we instantiate X with d0 and reify these levels w.r.t.\ their binders, we have
the problem solved with no surprise}: the two equations become \scm{(== (λ (a) a) (λ (b) b))} and
\scm{(== (λ (c) c) (λ (d) d))}. The paper goes on (p.\,3--4): \enquote{using grafting with de Bruijn
levels captures the essence of α-equivalence: names don't matter. A unification algorithm based on de
Bruijn levels would also be more straightforward because it's practically first-order. But, we cannot
yet lift the scope restriction on lambda variables and drop in de Bruijn levels, because sometimes
names do matter and \dots}. The sentence continues in the title of §3.

\begin{kanren}[(aux kanren micro) already grafts]
Grafting is what a first-order unifier does on nameless terms. Today's \micro{} unifies records slot
by slot, so a \scm{define-record} is a new term constructor for free. If we write λ-terms with de
Bruijn indices as records, \scm{=°} grafts. It reproduces the paper's answers, and also the problems
discussed below. The demonstration uses indices, not levels; for the examples here only the
outermost binder matters, and it is index~$0$ under one λ.
\begin{scheme}
(import scheme (chicken base) (chicken pretty-print)
        (aux base) (aux stream) (aux kanren micro))

(define-syntax-rule (show label expr)
  (begin (display ";; ") (display label) (newline)
         (pretty-print expr)))

;; 1. Named binders as plain lists: syntactic unification ignores alpha.
(show "named: (λ a a) = (λ b b)"
  (°->list #f (fresh° (q) (=° '(λ a a) '(λ b b)))))

;; 2. Nameless binders as records with de Bruijn indices.
(define-record lam body)
(define-record dbi n)
(define (λ. t) (make-lam t))
(define (ix n) (make-dbi n))

(show "de Bruijn: λ.0 = λ.0"
  (°->list #f (fresh° (q) (=° (λ. (ix 0)) (λ. (ix 0))))))
(show "grafting: λa.X = λb.b"
  (°->list #f (fresh° (x) (=° (λ. x) (λ. (ix 0))))))
(show "grafting: λa.X = λb.b  and  λc.X = λd.d"
  (°->list #f (fresh° (x) (=° (λ. x) (λ. (ix 0)))
                          (=° (λ. x) (λ. (ix 0))))))
(show "true = false: λa.λa.X = λc.λd.c"
  (°->list #f (fresh° (x) (=° (λ. (λ. x)) (λ. (λ. (ix 1)))))))
(show "dangling: λa.X = λb.b, then Q = X at top level"
  (°->list #f (fresh° (q) (fresh° (x) (=° (λ. x) (λ. (ix 0)))
                                      (=° q x)))))
\end{scheme}
\begin{answers}
;; named: (λ a a) = (λ b b)
()
;; de Bruijn: λ.0 = λ.0
((λ (α) α))
;; grafting: λa.X = λb.b
((λ () (make-record-instance 'dbi 0)))
;; grafting: λa.X = λb.b  and  λc.X = λd.d
((λ () (make-record-instance 'dbi 0)))
;; true = false: λa.λa.X = λc.λd.c
((λ () (make-record-instance 'dbi 1)))
;; dangling: λa.X = λb.b, then Q = X at top level
((λ () (make-record-instance 'dbi 0)))
\end{answers}
The first query shows that named binders as plain data are useless: α-equivalent terms do not
unify. The next three reproduce the paper's \texttt{'(\_0)}, \texttt{'(d0)}, \texttt{'(d0)}. The fifth
is the \enquote{true = false} problem of the next subsection. The answer \texttt{(dbi 1)} is
perfectly correct for the nameless terms, but it denotes the \emph{outer} binder. In the user's
named term $\abs{a}{\abs{a}{X}}$ that binder is shadowed. The last query shows a \emph{dangling}
index: $Q$ lives at top level, outside every binder, yet it is bound to \enquote{the nearest
binder}. \Cref{sec:lkanren-assess} returns to it.
\end{kanren}

\subsection{Binder-merging and \enquote{true = false}}\label{sec:lkanren-merge}

\begin{definition}[Binder-merging {\cite[p.\,4]{MaFriedman2021}}]\label{def:lkanren-merge}
\enquote{When two terms are unified, one term's binders are replaced by the other's. Then
α-equivalence is as simple as structural equivalence.} The paper says that it is \enquote{used in
all pattern unification algorithms, except for those with de Bruijn numbers}.
\end{definition}

\begin{example}[binder-merging at work, p.\,4]
\[
  \abs{x}{\abs{y}{X\,x\,y}} \lkq \abs{a}{\abs{b}{a\,b}}
  \;\longrightarrow\;
  \abs{a}{\abs{b}{X\,a\,b}} \lkq \abs{a}{\abs{b}{a\,b}}
\]
This is solved by $\{X\mapsto \abs{a}{\abs{b}{a\,b}}\}$. Instantiating gives
$\abs{a}{\abs{b}{((\abs{a}{\abs{b}{a\,b}})\,a\,b)}}$, which β-reduces to $\abs{a}{\abs{b}{a\,b}}$.
\end{example}

Now take a binder name that is reused, i.e.\ \emph{shadowing} (dashed box, p.\,4):

\begin{scheme}
> (run 1 X
    (== (λ (a) (λ (a) X)) (λ (c) (λ (d) c))))
\end{scheme}

\noindent \enquote{With either binder-merging or de Bruijn syntax, the unification problem is
effectively reduced to} $\abs{c}{\abs{d}{X}}\lkq\abs{c}{\abs{d}{c}}$. \enquote{Then it appears to be
solvable by $\{X\mapsto c\}$, or equivalently $\{X\mapsto \texttt{d0}\}$, that associates $X$ with the
outer-binder.} Naive reification, which writes each level back as the name of its binder, then gives

\begin{scheme}
(== (λ (a) (λ (a) a)) (λ (c) (λ (d) c)))
\end{scheme}

\noindent \enquote{which claims that the λ-encoded false is true!} With Church booleans
$\cst{true}=\abs{t}{\abs{f}{t}}$ and $\cst{false}=\abs{t}{\abs{f}{f}}$, the left-hand side
$\abs{a}{\abs{a}{a}}$ is α-equivalent to $\cst{false}$, because the occurrence $a$ refers to the inner
binder. The right-hand side is $\cst{true}$.

The paper argues that binder-merging goes wrong \enquote{even without lifting the scope
restrictions}:
\[
  \abs{x}{\abs{y}{x\,y}} \lkq \abs{a}{\abs{a}{a\,a}}
  \quad\text{is converted to}\quad
  \abs{a}{\abs{a}{a\,a}} \lkq \abs{a}{\abs{a}{a\,a}},
\]
\enquote{which then ends up with another false positive}. The correct answer is failure:
$\abs{a}{\abs{a}{a\,a}}\equiv_\alpha\abs{a}{\abs{b}{b\,b}}$ applies the \emph{inner} variable to itself.
The paper's claim: \enquote{binder-merging does not handle the name shadowing problem. This is
overlooked by all pattern unification algorithms. Algorithms that use what we call binder-merging state
upfront that they don't consider name shadowing}. Nipkow--Prehofer and Qian are cited as examples. The
existing remedies are \enquote{a pre-processor that adds gensyms to lambda variables to make each
individual unique or a smart reifier that detects name shadowing} (p.\,4). Their downsides are
unreadable gensyms and extra complexity. The related work gives a real case (p.\,8): Agda uses de Bruijn
numbers in its core language, so it \enquote{must make all variable names unique in a “gensym” manner},
which \enquote{bloats type checking with randomly generated symbols and makes the error messages painful to
read}.

\subsection{Static closures}\label{sec:lkanren-static}

\enquote{We propose a simple and efficient alternative, static closures. Static closures are
integrated into the unification process and they decide α-equivalence while correctly handling name
shadowing. We obtain a new algorithm by extending traditional pattern unification with static
closures} (p.\,5). No efficiency evidence or complexity bound is given.

\begin{definition}[Terms, scopes, static closures {\cite[Fig.~1, p.\,5]{MaFriedman2021}}]\label{def:lkanren-closure}
\[
\begin{array}{rcll}
t & ::= & X & \text{unification variable}\\
  & \mid & t\;t & \text{function application}\\
  & \mid & \abs{a}{t} & \text{named binder}\\
  & \mid & a & \text{name}\\
  & \mid & \lambda\,t & \text{nameless binder (derived)}\\
  & \mid & \ell & \text{level (derived)}\\[2pt]
\phi & ::= & \epsilon \mid a,\phi & \text{scope}\\
\kappa & ::= & \lkclo{\phi}{t} & \text{static closure}
\end{array}
\]
\enquote{The nameless syntax, i.e., nameless-binders and levels are derived during unification}
(p.\,5): the user writes only $X$, $t\,t$, $\abs{a}{t}$ and $a$. \enquote{A static closure is a pair
of a scope and a unification term. A scope is a list of binders. In essence, a static closure
captures the surrounding binders of a term. E.g., the term λa.λb.X reduces to the static closure
$\lkclo{a,b}{X}$.}
\end{definition}

Two remarks from our notes. First, the grammar has no constants, although the examples use them
(\texttt{'cat}); read \enquote{name} as covering free constants too. Second, the grammar conses on
the left ($a,\phi$), but the example writes the scope of $\abs{a}{\abs{b}{X}}$ as $a,b$, outermost
first. We adopt \emph{outermost first}: then the level of a binder is its position in the scope.

\begin{definition}[Operations on scopes {\cite[p.\,5]{MaFriedman2021}}]\label{def:lkanren-extdb}
In the paper's words:
\begin{itemize}
  \item \enquote{$\mathrm{ext}(\phi,a)$ adds the name $a$ to the scope $\phi$. If $a$ has been
    previously added to $\phi$, then this older one is shadowed and its corresponding level becomes
    unusable.}
  \item \enquote{$\mathrm{db}(\phi,a)$ is $\lkJust\;\ell$ where $\ell$ is the de Bruijn level of $a$
    w.r.t.\ the scope $\phi$, if $a$ is not shadowed in $\phi$; otherwise, the result is
    $\lkNothing$.}
  \item \enquote{$\mathrm{db}(\phi,\ell)$ is $\lkJust\;\ell$, if $\ell$ is available in $\phi$;
    otherwise, the result is $\lkNothing$.}
\end{itemize}
\end{definition}

So $\mathrm{db}$ is overloaded. On a name it converts the name to a level; on a level it checks that
the level is available. The notes make this precise. Write $\phi=a_0,\dots,a_{n-1}$ outermost first:
\[
  \mathrm{ext}(\phi,a)=\phi,a\qquad
  \mathrm{avail}(\phi)=\{\ell<|\phi| \;:\; a_\ell\notin\{a_{\ell+1},\dots,a_{n-1}\}\}
\]
\[
  \mathrm{db}(\phi,a)=\begin{cases}\lkJust\;\max\{\ell : a_\ell=a\} & a\in\phi\\ \lkNothing & a\notin\phi\end{cases}
  \qquad
  \mathrm{db}(\phi,\ell)=\begin{cases}\lkJust\;\ell & \ell\in\mathrm{avail}(\phi)\\ \lkNothing&\text{otherwise}\end{cases}
\]
In words: a level is \emph{available} if its binder's name is not rebound further in. The
name-to-level case needs an interpretation. An \emph{occurrence} of a name always refers to its
innermost binder, which is never shadowed, so the paper's \enquote{if $a$ is not shadowed} can only
mean \enquote{the innermost $a$}. The other reading would make $\abs{a}{\abs{a}{a}}$ impossible to
convert, which is absurd. This remains an interpretation.

\begin{tutorial}[the three examples of §4, step by step]
\textbf{Example 1} (p.\,5). Both sides are turned into closures:
\[
  \abs{a}{\abs{b}{X}} \lkq \abs{c}{\abs{d}{c}}
  \;\longrightarrow\; \lkclo{a,b}{X} \lkq \lkclo{c,d}{c}.
\]
\enquote{We first apply db on the right-hand side}: $\mathrm{db}((c,d),c)=\lkJust\;0$, so the problem
is $\lkclo{a,b}{X}\lkq\lkJust\;0$. \enquote{Then we try to solve X by applying db on the scope of the
left-hand side and 0}: $\mathrm{db}((a,b),0)=\lkJust\;0$, because $a$ is not rebound. Then
$\lkJust\;0\lkq\lkJust\;0$. \enquote{It works. So, the solution is $\{X\mapsto\texttt{d0}\}$.}

\medskip
\textbf{Example 2} (p.\,5--6), the shadowing case:
\[
  \abs{a}{\abs{a}{X}} \lkq \abs{c}{\abs{d}{c}}
  \;\longrightarrow\; \lkclo{a,a}{X} \lkq \lkclo{c,d}{c}.
\]
As before, the right side gives $\lkJust\;0$. \enquote{This time, db($\langle a, a; 0\rangle$) is
Nothing because the outer binder a is shadowed}. Here the paper puts the $0$ inside the closure
brackets, an abuse of notation for $\mathrm{db}((a,a),0)$. We get $\lkNothing\lkq\lkJust\;0$,
\enquote{and we know the unification fails}. True is not false any more.

\medskip
\textbf{Example 3} (p.\,6), two unknowns:
\[
  \abs{a}{\abs{b}{X}} \lkq \abs{c}{\abs{d}{Y}}
  \;\longrightarrow\; \lkclo{a,b}{X} \lkq \lkclo{c,d}{Y}.
\]
\enquote{Here, the grafting itself is not able to describe the solution. Now our solution needs a
context: the set of equations between static closures of unification variables. This problem has the
solution: the grafting being empty and the context being $\{\lkclo{a,b}{X}\lkq\lkclo{c,d}{Y}\}$.}
The paper adds that, in practice, \enquote{the accumulated equations between closures of one
unification problem are usually solved and reduced to grafting, during the unification on other
problems}.
\end{tutorial}

So a solution is a pair $(\theta,C)$: a grafting $\theta$ and a \emph{context} $C$ of postponed
flex-flex equations between closures. This is analogous to the flex-flex pairs that Huet's procedure
leaves unsolved (\cref{ch:hounif}).

\subsection{A reconstruction of the rules (not in the paper)}\label{sec:lkanren-recon}

To check that the three examples fit together, we wrote a minimal rule system that
reproduces all of them. It was implemented as a standalone CHICKEN program of about 150 lines, which
does not use \micro{}. We repeat that this is \textbf{our reconstruction, not the authors'
algorithm}.

\paragraph{Level form.} A surface term $t$ under scope $\phi$ is converted to internal syntax
$⟦t⟧_\phi$:
\[
  ⟦X⟧_\phi=\lkclo{\phi}{X}\qquad
  ⟦a⟧_\phi=\begin{cases}\ell & \mathrm{db}(\phi,a)=\lkJust\;\ell\\ a&\text{otherwise (a constant)}\end{cases}
\]
\[
  ⟦\abs{a}{t}⟧_\phi=\lambda\,⟦t⟧_{\mathrm{ext}(\phi,a)}\qquad
  ⟦t\,u⟧_\phi=⟦t⟧_\phi\,⟦u⟧_\phi
\]
Both sides of an equation are converted from $\epsilon$, so they are always compared at equal depth.
A binding is $\theta(X)=(\phi_0,t)$, where $t$ is in level form at depth $|\phi_0|$. Levels below
$|\phi_0|$ are \emph{free} (scope) levels; the others are bound inside $t$. To \emph{instantiate}
$\theta(X)$ at an occurrence $\lkclo{\phi}{X}$:
\begin{itemize}
  \item a free level $\ell$ must satisfy $\mathrm{db}(\phi,\ell)=\lkJust\;\ell$, or else the
    occurrence fails;
  \item a local level $\ell\ge|\phi_0|$ is shifted to $\ell+|\phi|-|\phi_0|$;
  \item a nested closure is re-rooted at $\phi$.
\end{itemize}

\paragraph{Rules.} The judgement is $\theta;C\vdash t\lkq u\Rightarrow\theta';C'$. Terms are walked
first, i.e.\ a bound closure is replaced by its instance. Any pair of head constructors not listed
fails.
\begin{mathpar}
\inferrule{ }{\theta;C\vdash \ell\lkq\ell \Rightarrow \theta;C}\;\rulename{Lvl}
\and
\inferrule{c\text{ a constant}}{\theta;C\vdash c\lkq c \Rightarrow \theta;C}\;\rulename{Const}
\and
\inferrule{\theta;C\vdash t\lkq u\Rightarrow\theta';C'}{\theta;C\vdash \lambda t\lkq\lambda u\Rightarrow\theta';C'}\;\rulename{Lam}
\and
\inferrule{\theta;C\vdash t_1\lkq u_1\Rightarrow\theta_1;C_1 \\ \theta_1;C_1\vdash t_2\lkq u_2\Rightarrow\theta_2;C_2}
          {\theta;C\vdash t_1\,t_2\lkq u_1\,u_2\Rightarrow\theta_2;C_2}\;\rulename{App}
\and
\inferrule{X\notin\mathrm{dom}\,\theta \\ FL_{<|\phi|}(t)\subseteq\mathrm{avail}(\phi) \\ X\notin FV(t) \\
           \theta[X\mapsto(\phi,t)];\emptyset\vdash C\Rightarrow\theta';C'}
          {\theta;C\vdash\lkclo{\phi}{X}\lkq t\Rightarrow\theta';C'}\;\rulename{Graft}
\and
\inferrule{ }{\theta;C\vdash\lkclo{\phi}{X}\lkq\lkclo{\phi}{X}\Rightarrow\theta;C}\;\rulename{Refl}
\and
\inferrule{X,Y\notin\mathrm{dom}\,\theta \\ (X,\phi)\neq(Y,\psi)}
          {\theta;C\vdash\lkclo{\phi}{X}\lkq\lkclo{\psi}{Y}\Rightarrow\theta;C\cup\{\lkclo{\phi}{X}\lkq\lkclo{\psi}{Y}\}}\;\rulename{Postpone}
\end{mathpar}
In \rulename{Graft}, $t$ is not an unbound closure. $FV(t)$ is the set of unification variables of
$t$ (the occurs check). $FL_{<|\phi|}(t)$ is the set of levels below $|\phi|$ that occur in $t$: those
that point into the scope $\phi$ rather than to a nameless binder inside $t$. The availability premise
says that a solution may mention only scope binders that are visible, i.e.\ not shadowed, at the
occurrence of $X$. The premise \enquote{$\vdash C$} means that every postponed equation is solved
again, the \emph{wake-up}. Example~2 fails in the availability premise of \rulename{Graft}:
$FL(0)=\{0\}\not\subseteq\mathrm{avail}(a,a)=\{1\}$.

\begin{tutorial}[a wake-up, step by step]
Take Example~3 and then bind $Y$, in two variants. Both start from
$\abs{a}{\abs{b}{X}}\lkq\abs{c}{\abs{d}{Y}}$ (first variant) or
$\abs{a}{\abs{a}{X}}\lkq\abs{c}{\abs{d}{Y}}$ (second variant).
\begin{enumerate}
  \item Conversion gives $\lambda\lambda\lkclo{a,b}{X}\lkq\lambda\lambda\lkclo{c,d}{Y}$. \rulename{Lam}
    twice, then \rulename{Postpone}: $\theta=\emptyset$, $C=\{\lkclo{a,b}{X}\lkq\lkclo{c,d}{Y}\}$. This is the
    line \texttt{(X ctx: 1)} below.
  \item The second equation $\abs{c}{\abs{d}{Y}}\lkq\abs{c}{\abs{d}{c}}$ converts to
    $\lambda\lambda\lkclo{c,d}{Y}\lkq\lambda\lambda\,0$. \rulename{Graft} checks
    $FL(0)=\{0\}\subseteq\mathrm{avail}(c,d)=\{0,1\}$ and binds $Y\mapsto((c,d),0)$.
  \item Wake-up: $C$ is solved again. Walking $\lkclo{c,d}{Y}$ instantiates $Y$'s solution at
    $(c,d)$, giving $0$. The equation becomes $\lkclo{a,b}{X}\lkq 0$, and \rulename{Graft} checks
    $0\in\mathrm{avail}(a,b)$. It holds, so $X\mapsto((a,b),0)$, printed \texttt{d0}.
  \item In the second variant the same wake-up yields $\lkclo{a,a}{X}\lkq 0$, but
    $\mathrm{avail}(a,a)=\{1\}$, so the whole conjunction fails.
\end{enumerate}
The second variant is \enquote{true = false} again, smuggled through the context. A wake-up that only
re-unified, without the availability check, would accept it.
\end{tutorial}

Running the prototype gives the following results. In \texttt{(λ d1)} the level is printed absolute,
at the depth where $X$ was solved.

\begin{answers}
X = cat                               => (cat ctx: 0)
λa.a = λb.b                           => (X ctx: 0)
λa.X = λb.b                           => (d0 ctx: 0)
λa.X = λb.b ∧ λc.X = λd.d             => (d0 ctx: 0)
λa.λa.X = λc.λd.c  (true=false)       => FAIL
λx.λy.(x y) = λa.λa.(a a)             => FAIL
Ex1 λa.λb.X = λc.λd.c                 => (d0 ctx: 0)
Ex2 λa.λa.X = λc.λd.c                 => FAIL
Ex3 λa.λb.X = λc.λd.Y                 => (X ctx: 1)
Ex3 then λc.λd.Y = λc.λd.c            => (d0 ctx: 0)
λa.λa.X = λc.λd.Y then Y = d0         => FAIL
λa.X = λb.λc.c ; λa.λb.X = λ_.λ_.λc.c => ((λ d1) ctx: 0)
λa.X = λb.λc.b                        => ((λ d0) ctx: 0)
λa.X = λb.b ∧ X = cat                 => FAIL
λa.X = λb.b ∧ λa.λa.X = λc.λd.d       => FAIL
λa.λb.X = λa.λb.a ∧ λb.λa.X = λc.λd.c => (d0 ctx: 0)
\end{answers}

\noindent Three lines are worth reading closely. In \enquote{λa.λa.X = λc.λd.Y then Y = d0}, the
wake-up must re-run the availability check, or the shadowing bug comes back through the context. In
\enquote{λa.X = λb.λc.c ; …}, a binder \emph{inside} a grafted solution must be shifted when the
solution is used at a different depth. The last line shows that the semantics is \emph{positional}:
the binder names above $X$ are irrelevant as long as no level is shadowed.

\subsection{The language: Horn clauses, \texttt{assume}, \texttt{all}}\label{sec:lkanren-mf21-lang}

§5 of the paper (p.\,6--7) recalls the logic from the 2020 paper \cite{MaLuFriedman2020}. Its Fig.~2
contrasts Horn clauses with hereditary Harrop formulas. In the figure the shared rows are
$D::=G\imp D\mid D\wedge D\mid\forall x\,D$ and $G::=G\wedge G\mid G\vee G\mid\exists x\,G$, and the
HH-only rows are $D\imp G$ and $\forall x\,G$ for goals. The atoms $A$ and $\top$ are missing from the
figure. With them it is exactly the book's grammar (3.1) \book{91}:
\[
  G ::= \top \mid A \mid G\wedge G \mid G\vee G \mid \exists x\,G \mid D\imp G \mid \forall x\,G\qquad
  D ::= A \mid G\imp D \mid D\wedge D \mid \forall x\,D .
\]
The four code boxes are these (all solid, p.\,6--7):

\begin{scheme}
> (run* (X Y)
    (conde
      [(== X Y) (== X 'cat)]
      [(== X 'cat) (== Y 'dog)]))
'((cat cat) (cat dog))

(define (empty? X)
  (== X '()))

>(run 2 X
   (assume (== 'apple 'orange)
     (== 'apple X)))
'(apple orange)

>(run 1 F
   (all (A)
     (== A (F A))))
'((lambda (_0) _0))
\end{scheme}

\noindent The paper prints \texttt{conde} with a superscript \textit{e}. The first box is
$\exists X\exists Y((X=Y\wedge X=\texttt{'cat})\vee(X=\texttt{'cat}\wedge Y=\texttt{'dog}))$. The
\texttt{define} is the clause $\forall X\,(X=\texttt{'()})\imp(\texttt{empty?}\;X)$. \texttt{assume}
\enquote{introduces hypotheses in its clause. As an example, we may use the assume operator to
temporarily extend the == relation}. Our reading, inferred because the paper gives no operational
account: the first answer comes from ordinary unification, and the second from the hypothesis used as
a clause whose head matches \scm{(== 'apple X)} with $X=\texttt{orange}$. For \texttt{all},
\enquote{we may synthesize a function that yields whatever it takes in}. Also, \enquote{the ==
operator recognizes β-equivalence. Indeed, it will be powered by the unification algorithm proposed
in the last section and recognize αβ-equivalence} (p.\,7). Note that \scm{(F A)} is a Miller pattern
(\cref{ch:hounif}): a variable applied to a distinct eigenvariable.

% =====================================================================================================
\section{A critical assessment of static closures}\label{sec:lkanren-assess}

The paper attacks pattern unification. The fair comparison also includes nominal unification
(\cref{sec:lkanren-nominal}), which the paper mentions only in related work. This section gives our own
analysis. In the λProlog column, $X$ (and $Y$) are quantified outside the λ, as in the paper's
queries.

\begin{table}[ht]
\centering\small
\caption{The paper's examples under three semantics (our analysis, after \cite{MaFriedman2021}).}
\label{tab:lkanren-three}
\begin{tabularx}{\textwidth}{>{\raggedright\arraybackslash}p{4.1cm}>{\raggedright\arraybackslash}X>{\raggedright\arraybackslash}X>{\raggedright\arraybackslash}X}
\toprule
problem & λProlog (patterns) & αKanren (nominal) & static closures\\
\midrule
$\abs{a}{a}\lkq\abs{b}{b}$ & \lkyes{} & \lkyes{} & \lkyes{}\\
$\abs{a}{X}\lkq\abs{b}{b}$ & \lkno{} (scope) & \lkyes{} $X=a$ & \lkyes{} $X=\texttt{d0}$\\
\dots\ $\wedge\ \abs{c}{X}\lkq\abs{d}{d}$ & \lkno{} & \lkno{} ($a\neq c$) & \lkyes{} \texttt{d0}\\
$\abs{a}{\abs{a}{X}}\lkq\abs{c}{\abs{d}{c}}$ & \lkno{} & \lkno{} (needs $a\#a$) & \lkno{}\\
$\abs{x}{\abs{y}{x\,y}}\lkq\abs{a}{\abs{a}{a\,a}}$ & \lkno{} & \lkno{} & \lkno{}\\
$\abs{a}{\abs{b}{X}}\lkq\abs{c}{\abs{d}{Y}}$ & \lkyes{} $X=Y$ (neither may depend on $a,b,c,d$) &
  $X=(a\,c)(b\,d)\cdot Y$ with $a\#Y$, $b\#Y$ & postponed\\
\bottomrule
\end{tabularx}
\end{table}

\begin{tutorial}[what λProlog offers the novice instead: make the dependency explicit]
The paper calls the failure of $\abs{a}{X}\lkq\abs{b}{b}$ counter-intuitive. Miller's 1991 paper
\cite[p.\,4--5]{Miller1991} explains why it is the intended behaviour, and how to get the novice's
answer anyway. Its program is the usual \texttt{append}, and its goal is
$\forall y\,(\mathtt{append}\;(\mathtt{cons}\;a\;(\mathtt{cons}\;b\;\mathtt{nil}))\;y\;X)$.
\rulename{generic} replaces $y$ by a new constant $k$, and backchaining twice would need
$X=\mathtt{cons}\;a\;(\mathtt{cons}\;b\;k)$. This is rejected, and Miller says why: \enquote{Such a failure
here is quite sensible since the value of $X$ should be independent of the choice of the constant used
to instantiate $\forall y$.} If we do want a value that depends on $y$, we must say so, by writing
$(H\;y)$ instead of $X$. Then $(H\;k)$ must equal $\mathtt{cons}\;a\;(\mathtt{cons}\;b\;k)$. Two λ-terms
would do, $\abs{w}{\mathtt{cons}\;a\;(\mathtt{cons}\;b\;k)}$ and
$\abs{w}{\mathtt{cons}\;a\;(\mathtt{cons}\;b\;w)}$, but \enquote{since $H$ cannot contain $k$ free, only
the second} is legal. Miller calls this \enquote{discharging the constant $k$}.

The same move answers the novice. Write $\abs{a}{(X\;a)}\lkq\abs{b}{b}$: this is a pattern ($X$ applied
to one bound variable), and its most general unifier is $X\mapsto\abs{x}{x}$. The conjunction with
$\abs{c}{(X\;c)}\lkq\abs{d}{d}$ is solved by the same $X$, and a top-level $Q=X$ gets the closed term
$\abs{x}{x}$, with nothing dangling. Both λKanren papers show this style of query:
\scm{(== (tie (a b) (app a b)) (tie (x y) (app f x y)))} of 2020 and
\scm{(all (A) (== A (F A)))} of 2021 (\cref{sec:lkanren-mf21-lang,sec:lkanren-mk20}). Static closures
make this dependency \emph{implicit}: $\lkclo{\phi}{X}$ behaves like $X$ applied to the available
levels of $\phi$ (observation~1 below).
\end{tutorial}

\subsection{Three counter-examples}\label{sec:lkanren-counter}

The notes isolate three examples that pin down what static closures really are. They are worth
keeping as regression tests for any implementation, whatever it chooses.

\begin{tutorial}[Counter-example 1: the conjunction separates static closures from nominal unification]
Problem: $\abs{a}{X}\lkq\abs{b}{b}\ \wedge\ \abs{c}{X}\lkq\abs{d}{d}$.

\emph{Nominal}, using the rule $a.t\approx^? b.s\Rightarrow t\approx^?(a\,b)\cdot s,\ a\#^? s$ of
\cref{sec:lkanren-nominal}. The first equation gives $X\approx^?(a\,b)\cdot b=a$ with $a\#b$, which
holds, so $X:=a$. The second becomes $c.a\approx^?d.d$, i.e.\ $a\approx^?(c\,d)\cdot d=c$, which
fails. Nominal unification thus agrees with static closures on the single equation ($X$ is the bound
name) and disagrees on the conjunction.

\emph{Pattern}: both equations fail at once, by the scope restriction.

\emph{Static closures}: $X\mapsto\texttt{d0}$ satisfies both.

So static closures are \emph{neither} pattern unification \emph{nor} nominal unification. $X$ denotes
an open term \emph{relative to the binder depth}. This is grafting on de Bruijn terms \emph{without}
the \enquote{pre-cooking} step of Dowek et al.\ (\cref{sec:lkanren-dhkp}).
\end{tutorial}

\begin{tutorial}[Counter-example 2: the same variable at two depths dangles]
Problem: \scm{(== Q X)} together with \scm{(== (λ (a) X) (λ (b) b))}.

Under static closures, $X\mapsto\texttt{d0}$, so the top-level answer is $Q=\texttt{d0}$. That is a
level which is \emph{not available} in the empty scope where $Q$ lives. The answer is meaningful only
relative to the occurrence of $X$ under one binder. The reading becomes incoherent as soon as the same
variable occurs at different depths. Our \micro{} run above shows the same thing with indices. The
prototype chooses to check availability inside unification, so \enquote{λa.X = λb.b then top-level
Y = X} fails there. But it skips the check when printing, which is how the paper can print
\texttt{'(d0)}. Dowek et al.\ fix this with $X[\uparrow^n]$; Teyjus and ELPI fix it with levels on
variables. The 2021 paper does not address it. It also gives no rules for application and β, nor for
flex-flex closures beyond postponing them.
\end{tutorial}

\begin{tutorial}[Counter-example 3: \enquote{true = false} is a bug of binder-merging, not of pattern unification]
Problem: $\abs{a}{\abs{a}{X}}\lkq\abs{c}{\abs{d}{c}}$.

\emph{Pattern}: $X$ cannot depend on the bound variables at all, so this fails.

\emph{Nominal}: $a.(a.X)\approx^? c.(d.c)$ gives $a.X\approx^?(a\,c)\cdot(d.c)=d.a$, with
$a\#(d.c)$, which holds. Then $a.X\approx^?d.a$ gives $X\approx^?(a\,d)\cdot a=d$, with
\textbf{$a\#a$}, which is false. So it fails.

\emph{Static closures}: fails (Example~2).

All three agree. The \enquote{false = true} danger is specific to \emph{binder-merging}
implementations that do not rename apart. The earlier λKanren paper renames apart: its ξ step uses
gensyms (\cref{sec:lkanren-mk20}), so it does not have the bug.
\end{tutorial}

\subsection{Further observations}

The notes add several observations that a porter must know. The paper does not say them.

\begin{enumerate}
  \item \textbf{The meaning of a unification variable changes.} In λProlog, $X$ in $\abs{a}{X}$
    cannot mention $a$. With static closures, $X$ implicitly depends on the \emph{positions} of its
    enclosing binders. Logically, $\lkclo{\phi}{X}$ behaves like Miller's raising
    $X\,\ell_{i_1}\cdots\ell_{i_k}$ over the available levels of $\phi$, with one twist: the
    \enquote{function} is indexed by absolute position, so occurrences under scopes of different
    lengths share it. The paper's related work quotes Huet's \enquote{elementary substitution}:
    \enquote{Each unification variable should capture all its visible bound variables} (p.\,8).
    Static closures make that idea implicit, minus the shadowed binders.
  \item \textbf{α-equivalence is no longer a congruence on open terms.} $\abs{a}{\abs{a}{X}}$ and
    $\abs{a}{\abs{b}{X}}$ are α-equivalent in the usual sense, because $X$ is a metavariable with no
    free names. Yet against the same $\abs{c}{\abs{d}{c}}$, Example~2 fails and Example~1 succeeds.
    With grafting, the names of the binders above a metavariable matter. Nominal logic manages exactly
    this with suspended permutations $\pi\cdot X$ and freshness $a\#X$. The paper's \enquote{names
    don't matter} holds for the levels, not for the closures.
  \item \textbf{Levels are stable for scope variables, not for binders inside solutions.} A grafted
    $\lambda.\ell_1$ solved at depth~1 must become $\lambda.\ell_2$ at depth~2. A locally nameless
    representation avoids the shift \cite{McBrideMcKinna2004,Chargueraud2012}: levels for scope
    variables, indices for binders inside solutions.
  \item \textbf{Flex-flex need not be postponed.} With levels, $\lkclo{\phi}{X}\lkq\lkclo{\psi}{Y}$
    at equal depth has a most general solution: $X\mapsto Y$, plus the restriction that $Y$ may use
    only levels in $\mathrm{avail}(\phi)\cap\mathrm{avail}(\psi)$. This is Miller-style pruning,
    phrased as a freshness-like constraint.
  \item \textbf{The trick must not reach \texttt{all}.} For ∀-goals the eigenvariable must respect the
    scope restriction, or \scm{(run 1 X (all (A) (== X A)))} would succeed. That is unsound for ∀.
    The level semantics is meant only for λ-binders inside terms.
  \item \textbf{β brings undecidability back.} Once \scm{==} \enquote{recognizes αβ-equivalence},
    unification outside the pattern fragment is undecidable (\cref{ch:hounif}). One must restrict to
    patterns, delay, or enumerate.
\end{enumerate}

The paper's related work (p.\,8) recalls that \enquote{nominal unification and pattern unification are
equally expressive}. Cheney reduces pattern unification to nominal unification \cite{Cheney2005}, and
Levy and Villaret give the converse reduction \cite{LevyVillaret2008}. Pattern unification is linear
(Qian, although the paper notes that Qian's algorithm \enquote{has not been implemented}); for nominal
unification \enquote{$O(n^2)$ is the state-of-the-art} (p.\,8). The notes summarize the correspondence like this:
the static closure $\lkclo{\phi}{X}$ is close to the nominal suspension $\pi\cdot X$ (both record the
context around an unknown); $\mathrm{db}$ is close to the action of $\pi$ on names; and \enquote{level
unavailable in $\phi$} is close to the freshness $\ell\#X$. But static closures are
\emph{positional} where nominal terms are \emph{name-based}, as counter-example~1 shows. The design is
closer to de Bruijn-level \enquote{contextual} metavariables than to nominal logic.

\begin{kanren}[what static closures would demand of a miniKanren]
The notes list what a static-closure layer on \micro{} would need. This is a requirements list, not a
design (see \cref{ch:design}):
\begin{itemize}
  \item \textbf{New term kinds, dispatched early.} λ, level and closure records must be handled
    \emph{before} the generic record branch. Today \scm{μkanren-structure?} matches any record
    instance and unifies it slot by slot, names included. The same applies to every walker: the occurs
    check, \scm{μkanren-state-find/value}, reification, \scm{μkanren-hasvar?}, and so on.
  \item \textbf{Conversion at the boundary.} Converting both sides of \scm{=°} to level form under
    $\epsilon$ leaves a unifier that is \enquote{practically first-order} except for closures.
  \item \textbf{A walk that can fail.} Walking $\lkclo{\phi}{X}$ instantiates $X$ at $\phi$, which can
    fail on an unavailable level. So the walk needs a third outcome besides \enquote{value} and
    \enquote{unbound}. The first-order hot path should stay untouched.
  \item \textbf{A context store with wake-up.} The postponed equations must be re-checked when their
    variables get bound. The \scm{μkanren-verify-D/touched} idiom, which re-runs only the entries that
    mention a just-bound variable, already does this for disequalities.
  \item \textbf{Disequality.} With postponement, a tentative unification inside \scm{≠°} may succeed
    by only adding a context equation, producing no bindings. An empty prefix in D means
    \enquote{already equal}, so this would be wrong. A disequality entry would then have to carry
    its postponed part as well.
\end{itemize}
\end{kanren}

% =====================================================================================================
\section{λKanren 2020: Miller's algorithm in a shallow embedding}\label{sec:lkanren-mk20}

The earlier paper, \emph{Higher-order Logic Programming with λKanren} \cite{MaLuFriedman2020}, is the
\enquote{[17]} of the 2021 paper. Its workshop title is \enquote{λKanren: Higher-order Logic
Programming with Shallow Embedding}. It adds four operators to μKanren \cite{HemannFriedman2013}:
\scm{tie} (a binder), \scm{app} (application), \scm{assume-rel} (an implication goal $D\imp G$) and
\scm{all} (a universal goal $\forall x\,G$). In it, \scm{==} \enquote{recognizes αβ-conversions}. The
paper claims \enquote{about 40 lines} added to μKanren (p.\,7); the repository's \texttt{unify.rkt}
has 648 lines.

\paragraph{Terms and state.} The terms are
$\mathit{Term}::=\mathit{Constant}\mid\mathit{Var}\mid\mathsf{App}(\mathit{rator},\mathit{rand})\mid
\mathsf{Tie}(\mathit{var},\mathit{body})\mid(\mathit{Term}\,.\,\mathit{Term})$. Pairs stay a separate
first-order constructor. The state is the list \texttt{(qs es σ name->idx)}. Here \texttt{qs} is an
\textbf{explicit quantifier prefix} of \texttt{(∃ x)} and \texttt{(∀ x)} entries, appended at the end
by \texttt{fresh} and \texttt{all}. \texttt{es} holds postponed equations and σ is the substitution.

\paragraph{The unifier} is a direct transcription of Miller 1991 §6 \cite{Miller1991}: ξ, rigid-rigid,
raising, pruning, flex-flex and flex-rigid (\cref{sec:lkanren-miller}). Its ξ step strips the
\texttt{Tie}s on both sides and renames both binder lists to the same \emph{fresh} gensyms by
capture-avoiding substitution. It then pushes a \texttt{(∀ w)} onto the prefix for each name. So
shadowing is handled by renaming apart, and the 2021 problem does not arise. Non-pattern applications
are \textbf{not delayed}; they make the unifier fail.

\begin{scheme}
(define-syntax tie
  (syntax-rules ()
    [(_ () t) t]
    [(_ (x0 x ...) body)
     (let ([x0 (Var 'x0)])
       (Tie x0 (tie (x ...) body)))]))

(define-syntax app
  (syntax-rules ()
    [(_ rator rand) (App rator rand)]
    [(_ rator rand0 rand ...)
     (app (App rator rand0) rand ...)]))
\end{scheme}

\noindent (The paper prints \texttt{[(\_ () t) body]}; the repository has \texttt{t}.) The paper reports
these unification queries (p.\,2--3):

\begin{scheme}
(run 1 q (== (tie (a b) (app a b)) (tie (x y) (app x y))))      ; ⇒ '(_0)
(run 2 f (== (tie (a b) (app a b)) (tie (x y) (app f x y))))    ; ⇒ '((tie (_0) (tie (_1) (app _0 _1))))
(run 1 f (== (tie (a) a) (tie (b) (app f b b))))                ; ⇒ '()   not a pattern
(run 1 q (== (tie (a b) (app a b)) (tie (x y) (app x q))))      ; ⇒ '()   scope restriction
\end{scheme}

\noindent The second query is a pattern, so \scm{run 2} still finds one answer, the most general
unifier: \enquote{There is only one instantiation} (p.\,3). Binders and unification variables are both
reified as \texttt{\_}$n$ symbols. The third is not a pattern ($f$ is applied to $b$ twice) and fails,
even though it has solutions ($\abs{x}{\abs{y}{x}}$ and $\abs{x}{\abs{y}{y}}$); in the paper's words,
\enquote{we cannot decide which b takes control}. The fourth fails by the scope restriction, which the
paper phrases as visibility: \enquote{A binding variable x is visible to a unification variable q if the
introduction of x lexically precedes that of q} (p.\,3). Here $q$ would have to be the bound $y$, which
is introduced after $q$.

\paragraph{Implication and the shallow-embedding problem.} The motivating program is a curriculum:

\begin{scheme}
(defrel (taken name class)
  (conde
    [(== 'Josh name) (== 'B521 class)]
    [(== 'Josh name) (== 'B522 class)]))
(defrel (pl-major name)
  (taken name 'B521)
  (taken name 'B523)
  (taken name 'B522))
> (run 1 q
    (assume-rel [(taken name class) (== 'Josh name) (== q class)]
      (pl-major 'Josh)))
'(B523)
\end{scheme}

\noindent The hypothetical clause extends \texttt{taken}, and \texttt{pl-major}, which was defined
\emph{before}, must see it. In a shallow embedding, relation definitions live in the host language's
lexical scope, so \texttt{pl-major} cannot see the new clause. The paper describes a
\textbf{dynamic-scope layer}: a map \texttt{name->idx} threaded through the query, a global
\texttt{idx->def}, and calls that go through both. (Our notes write \scm{(== q class)}; the extracted
text has a quote, \texttt{'q}, which could not produce \texttt{B523}.) Concretely, before the query
\texttt{name->idx} is \texttt{((pl-major . 2) (taken . 1))}; \texttt{assume-rel} pushes a new index, so
inside it the map is \texttt{((taken . 3) (pl-major . 2) (taken . 1))}, and the newer pair shadows the
older one when \texttt{pl-major} calls \texttt{taken} (p.\,5).

The paper also offers \texttt{assume}, \enquote{a shorter version of the assume-rel operator} that
\enquote{only takes an atomic hypothesis} (p.\,5), and uses it on a relation \texttt{eq} that is
reflexive, transitive and symmetric. Plain \scm{(run 1 q (eq 'apple 'orange))} does not terminate in
μKanren, because the recursive clauses never run out (p.\,5). With hypotheses (p.\,5--6):

\begin{scheme}
(defrel (eq x y)
  (conde
    [(== x y)]
    [(fresh (z) (eq x z) (eq z y))]
    [(eq y x)]))
> (run 5 q (assume (eq 'orange 'apple) (assume (eq 'orange 'dog) (eq 'orange q))))
'(dog apple orange orange dog)
> (run 1 q (assume (eq 'orange q) (eq 'apple 'orange)))
'(apple)
\end{scheme}

\noindent In the first query the answers come from the two hypotheses (\texttt{dog}, \texttt{apple}) and
from reflexivity (\texttt{orange}); \texttt{run 5} reports repetitions, because the recursive clauses
derive the same facts again along different paths. The second query \enquote{runs backward}: \enquote{the
hypothesis can be inferred as well} (p.\,6). The hypothesis \scm{(eq 'orange q)} mentions the query
variable $q$, and it is \emph{not} renamed when it is used. The goal \scm{(eq 'apple 'orange)} can use the
symmetry clause to become \scm{(eq 'orange 'apple)}, which matches the hypothesis with $q=\texttt{apple}$.
This is the book's semantics for a clause with a free logic variable (\cref{ch:hh}): only a clause's own
∀-bound variables are freshened per use.

\paragraph{Universal goals.} \enquote{the all operator effectively creates a constant \dots\ we create
an all variable as a free binding variable. Thus, the all variable cannot be unified with anything
but itself} (p.\,6). The scope restriction is enforced by the prefix. The examples (p.\,6--7):

\begin{scheme}
> (run 1 q (all (x)
             (assume-rel (taken x 'B521)
               (assume-rel (taken x 'B522)
                 (assume-rel (taken x 'B523)
                   (pl-major x))))))
'(_0)
> (run 1 f (all (x) (== x (app f x))))
'((tie (_0) _0))
(defrel (proved x) (== x 'trivial))
> (run 1 g (all (p) (proved p)))            '()
> (run 1 g (proved g))                      '(trivial)
> (run 1 g (all (p q) (assume ((proved p) (proved q)) (proved q) (proved p))))   '(_0)
> (run 1 goal (all (p q) (assume ((proved p)) (conde [(proved p)] [(proved q)]))))  '(_0)
\end{scheme}

\noindent The query \scm{(all (p) (proved p))} shows \rulename{generic} at work. \texttt{proved} holds only of
\texttt{trivial}, and a generic $p$ is not \texttt{trivial}, so $\forall p.\,\texttt{proved}\;p$ is not
provable, whereas the same goal without \scm{all} succeeds with $g=\texttt{trivial}$. The last two queries
use λKanren \enquote{as a theorem prover} (§5, p.\,6--7). The first proves the commutativity of
conjunction, $\forall p,q\,(p\wedge q)\imp(q\wedge p)$: \rulename{generic} makes $p$ and $q$ new
constants, \rulename{augment} adds \scm{(proved p)} and \scm{(proved q)} as hypotheses, and the two
conjuncts of the body are then proved by backchaining on them. The second proves disjunction introduction,
$\forall p,q\;p\imp(p\vee q)$, by the first \scm{conde} branch. The answers \texttt{'(\_0)} only say that
the query variable was left unconstrained: the information is in the success.

\begin{pitfall}[hypotheses in a host dynamic parameter get lost]
The repository's active code does not use the described state-threaded design. It uses Racket
parameters: \texttt{defrel} makes each body a \texttt{make-parameter}, and \texttt{assume} is a
\texttt{parameterize} that extends it. We reproduced this structure in CHICKEN. The body
of \texttt{pl-major} is only \emph{constructed} when the goal is applied to a state, which happens
after \texttt{parameterize} has returned. So the inner call to \texttt{taken} reads the
\emph{original} parameter, and the output was \texttt{(orig Josh B523)}. Relation calls nested inside
other relations' delayed bodies therefore lose the hypotheses. The described \texttt{name->idx} design
does not have this problem.
\end{pitfall}

\begin{kanren}[hypotheses belong to the state]
\micro{}'s \scm{define-relation} is inverse-η-delayed: it expands to
\scm{(λ (s) (δ ((and° g ...) s)))}. This is exactly the situation of the experiment above. A
hypothetical clause database for \rulename{augment} must therefore be a \emph{field of the state},
pushed on entry to $G$ and invisible outside it. It can never be a host parameter. A further remark
from the notes: a hypothetical clause $\forall\bar x(G\imp A)$ can be a host closure whose body runs
\scm{fresh°}. Each use then renames the clause's own variables, while the variables captured from
the enclosing scope stay shared. This is exactly λProlog's semantics for clauses with free logic
variables. How the module does this is shown in \cref{sec:module-hh}.
\end{kanren}

% =====================================================================================================
\section{Nominal unification and αKanren}\label{sec:lkanren-nominal}

Nominal unification \cite{UrbanPittsGabbay2004} takes a different route to α-equivalence. Bound names
stay names, called \emph{atoms}. Metavariables are holes, and substitution of a term for a
metavariable is \emph{capturing}, i.e.\ grafting. Capture is controlled by two devices:
\emph{permutations} suspended on variables, and \emph{freshness} constraints. Our presentation of the
rules follows Urban's survey \cite{Urban2010}.

\begin{definition}[Nominal terms and permutations]\label{def:lkanren-nominal}
\[
  t ::= \langle\rangle \mid \langle t_1,t_2\rangle \mid f\,t \mid a \mid a.t \mid \pi\cdot X
\]
where $a.t$ abstracts the atom $a$ in $t$, and $\pi\cdot X$ is a variable with a suspended
permutation. A permutation $\pi$ is a list of swappings $(a\,b)$:
\[
  []\cdot a=a,\qquad ((a_1\,a_2)::\pi)\cdot a=\begin{cases}a_2 & \pi\cdot a=a_1\\ a_1 & \pi\cdot a=a_2\\ \pi\cdot a&\text{otherwise.}\end{cases}
\]
Composition is list append and the inverse is reversal. The \emph{disagreement set} is
$ds(\pi,\pi')=\{a\mid \pi\cdot a\neq\pi'\cdot a\}$. A permutation acts on an abstraction on
\emph{both} the binder and the body, and it suspends on variables:
$\pi\cdot(\pi'\cdot X)=(\pi @ \pi')\cdot X$. Substitution is capturing: $\sigma(a.t)=a.\sigma(t)$ with
no side condition, and $\sigma(\pi\cdot X)=\pi\cdot\sigma(X)$.
\end{definition}

\begin{definition}[Freshness and α-equality {\cite[Fig.~1]{Urban2010}}]\label{def:lkanren-fresh}
With a freshness context $\nabla$, a set of pairs $a\#X$:
\begin{mathpar}
\inferrule{ }{\nabla\vdash a\#\langle\rangle}
\and
\inferrule{\nabla\vdash a\#t_1 \\ \nabla\vdash a\#t_2}{\nabla\vdash a\#\langle t_1,t_2\rangle}
\and
\inferrule{\nabla\vdash a\#t}{\nabla\vdash a\#f\,t}
\and
\inferrule{ }{\nabla\vdash a\#a.t}
\and
\inferrule{\nabla\vdash a\#t \\ a\neq b}{\nabla\vdash a\#b.t}
\and
\inferrule{a\neq b}{\nabla\vdash a\#b}
\and
\inferrule{(\pi^{-1}\cdot a,X)\in\nabla}{\nabla\vdash a\#\pi\cdot X}
\\
\inferrule{\nabla\vdash t_1\approx t_2}{\nabla\vdash a.t_1\approx a.t_2}
\and
\inferrule{a\neq b \\ \nabla\vdash a\# t_2 \\ \nabla\vdash t_1\approx(a\,b)\cdot t_2}{\nabla\vdash a.t_1\approx b.t_2}
\and
\inferrule{\forall c\in ds(\pi,\pi').\ (c,X)\in\nabla}{\nabla\vdash\pi\cdot X\approx\pi'\cdot X}
\end{mathpar}
plus the congruence rules for unit, pairs, function symbols and atoms.
\end{definition}

The intuition: $a\#t$ says that $a$ does not occur \emph{free} in $t$, also after any later
instantiation of the variables of $t$. To compare $a.t_1$ with $b.t_2$, rename $b$ to $a$ in $t_2$ by
a \emph{swapping}, which is always safe because it is a bijection, and require that $a$ was not
already free in $t_2$.

\begin{definition}[Nominal unification {\cite[p.\,9]{Urban2010}}]
\[
  \{\langle t_1,t_2\rangle\approx^?\langle s_1,s_2\rangle,\dots\}\Rightarrow\{t_1\approx^?s_1,\ t_2\approx^?s_2,\dots\}\qquad
  \{a.t\approx^?a.s,\dots\}\Rightarrow\{t\approx^?s,\dots\}
\]
\[
  \{a.t\approx^?b.s,\dots\}\Rightarrow\{t\approx^?(a\,b)\cdot s,\ a\#^?s,\dots\}\qquad
  \{\pi\cdot X\approx^?\pi'\cdot X,\dots\}\Rightarrow\{a\#^?X\mid a\in ds(\pi,\pi')\}\cup\{\dots\}
\]
and $\pi\cdot X\approx^?t$, with $X$ not in $t$, gives $X:=\pi^{-1}\cdot t$. A solution is a pair
$(\nabla,\sigma)$ with $\nabla\vdash\sigma(t_1)\approx\sigma(t_2)$. For example, $a.X\approx^?b.X$ has
the solution $(\{a\#X,\,b\#X\},[\,])$.
\end{definition}

Most general solutions exist, and \textbf{no fresh names are ever generated}: every atom of a solution
occurs in the problem. Naive rules are exponential; quadratic algorithms exist.

\subsection{αKanren}

αKanren \cite{ByrdFriedman2007,Byrd2009} embeds nominal unification in miniKanren. \scm{fresh}
introduces \emph{noms}: \enquote{a nom behaves more like a constant than a variable, since it only
unifies with itself or with an unassociated variable}. The constructor \scm{(tie a t)}, typeset
$\bowtie$, binds the nom $a$ in $t$; it delimits scope and does not create noms. \scm{(# a t)},
entered \scm{hash}, asserts that $a$ does not occur free in $t$. Ordinary logic variables are
introduced by \scm{exist}. The thesis examples (ch.~9; \texttt{◃▹} is the thesis glyph for tie):

\begin{scheme}
(run* (q) (fresh (a b) (≡ (◃▹ a a) (◃▹ b b))))      ⇒ (_.0)
(run* (q) (fresh (a b) (≡ (◃▹ a q) (◃▹ b b))))      ⇒ (a.0)
(run* (q) (fresh (a b) (# a (◃▹ b a))))             ⇒ ()
(run* (q) (fresh (a) (# a (◃▹ a a))))               ⇒ (_.0)
(run* (q) (fresh (a b) (exist (x y) (≡ (◃▹ a (◃▹ a x)) (◃▹ a (◃▹ b y))) (≡ `(,x ,y) q))))
  ⇒ ((((susp ((a.0 a.1)) _.0) _.0) : ((a.0 . _.0))))
\end{scheme}

\noindent The second line is counter-example~1's single equation: $q$ becomes the bound nom. The last
line shows a suspension in an answer, with a freshness constraint printed after the colon. The
capture-avoiding substitution relation shows the programming style:

\begin{scheme}
(define substo
  (λ (e new a out)
    (matche (e out)
      (((var ,a) ,new))
      (((var ,y) (var ,y)) (# a y))
      (((app ,rator ,rand) (app ,rator-res ,rand-res))
       (substo rator new a rator-res)
       (substo rand new a rand-res))
      (((lam ,(◃▹ @c body)) (lam ,(◃▹ c body-res)))
       (# c a)
       (# c new)
       (substo body new a body-res)))))
(run* (q) (fresh (a b) (substo `(lam ,(◃▹ a `(app (var ,a) (var ,b)))) `(var ,b) a q)))
  ⇒ ((lam (tie a.0 (app (var a.0) (var a.1)))))     ; [b/a]λa.ab = λc.cb
\end{scheme}

\noindent (\texttt{@c} in a \texttt{matche} pattern introduces $c$ as a nom.) The unifier works in two
phases. First the σ-rules turn equations into a substitution plus pending freshness constraints. Then
the ∇-rules reduce freshness constraints until each is $a\#x$ with $x$ an unbound variable. Variables
are suspensions with an empty permutation. A triangular-substitution version, due to Joe Near, is
reported \enquote{as much as ten times faster}. The thesis states a limitation (p.\,178): \enquote{neither
currently supports disequality constraints}. αKanren also powers αleanTAP \cite{NearByrdFriedman2008},
where ∀-bound variables of the object logic are noms. There, \enquote{Universally quantified formulas
are used as templates, rather than instantiated directly}.

\begin{kanren}[what nominal unification asks of a miniKanren]
It needs three new term kinds (nom, \scm{tie}, suspension $\pi\cdot X$) and a freshness store ∇ of
$a\#X$ pairs, re-reduced after each unification. That re-reduction works like \micro{}'s
post-\scm{=°} verification of D, A and T. Note that \scm{absent°} is \emph{not} freshness: freshness is
about \emph{free} occurrences, so \scm{(# a (◃▹ a a))} succeeds, whereas absento forbids any
occurrence. Nominal unification has no function variables and no β, so it is a different logic from
λProlog's. Noms are distinct constants, not eigenvariables, and ∀-goals still need their own
mechanism (see \cref{ch:design}).
\end{kanren}

% =====================================================================================================
\section{Eigenvariables, negation and ∀ inside miniKanren}\label{sec:lkanren-eigen}

\subsection{The scope-extrusion problem}

\rulename{generic} proves $\forall x.G$ by proving $G[c/x]$ for a \emph{new} constant $c$. If the
constant is only a gensym, nothing stops an older logic variable from being bound to it later. That
\enquote{proves} $\exists X\forall y.\,X=y$, which must fail. Contrast $\forall y\exists X.\,X=y$,
which must succeed. The fixes in the literature are exactly four:
(i)~an explicit prefix \cite{Miller1991,MaLuFriedman2020}; (ii)~universe indices or levels \cite{NadathurJayaramanKwon1995,DunchevEtAl2015}; (iii)~an absento-like eigen constraint over the variables in scope \cite{HemannByrdFriedman2013}; (iv)~an end-of-scope check.
The notes call (iv) the weakest. It is sound only if nothing created inside the scope survives it in a
form that can be bound later, and its failures are detected late.

\subsection{\texttt{eigen} in miniKanren with symbolic constraints}

Hemann, Byrd and Friedman added \scm{eigen} to miniKanren-with-symbolic-constraints in 2013
\cite{HemannByrdFriedman2013}. The header comment of \texttt{mk.scm} reads: \enquote{E = (e* . x*)*,
where e* is a list of eigens and x* is a list of variables. Each e in e* is checked for any of its
eigens be in any of its x*. Then it fails.} The state carries B, the list of all variables created so
far. \scm{eigen} creates fresh eigen objects and installs an eigen-absento constraint $(e^*\,.\,B)$:
\emph{none of the variables that existed before the eigen may ever walk to a term containing it}.

\begin{scheme}
(define-syntax eigen
  (syntax-rules ()
    ((_ (x ...) g0 g ...)
     (lambdag@ (c : B E S)
       (let ((x (eigen-var)) ...)
         ((fresh () (eigen-absento `(,x ...) B) g0 g ...) c))))))

(define eigen-absento
  (lambda (e* x*)
    (lambdag@ (c : B E S D Y N T)
      (cond
        ((eigen-occurs-check e* x* S) (mzero))
        (else (unit `(,B ((,e* . ,x*) . ,E) ,S ,D ,Y ,N ,T)))))))
\end{scheme}

\noindent This is the universe-index check implemented with explicit sets instead of numbers. Its
cost is $O(|B|\cdot|\text{term}|)$ per check. The README warns: \enquote{Beware: this implementation
does \emph{not} support use of \texttt{eigen} with constraints other than \texttt{==}.}

\begin{kanren}[the eigen-absento idea, by hand, on today's engine]
\micro{} already has \scm{absent°}, which takes any datum as its tag, including a record. So the
Hemann--Byrd--Friedman check can be written by hand. Here there is one outer variable $Q$, and the
eigen is a fresh record that $Q$ may never contain. This illustrates the idea; it is not an API
proposal.
\begin{scheme}
(import scheme (chicken base) (chicken pretty-print)
        (aux base) (aux stream) (aux kanren micro))

(define-syntax-rule (show label expr)
  (begin (display ";; ") (display label) (newline)
         (pretty-print expr)))

;; A generic constant as a fresh, unique record: no scope check at all.
(define-record eigen id)
(show "naive: exists Q forall x. Q = x"
  (°->list #f (fresh° (q) (=° q (make-eigen 1)))))

;; The eigen-absento idea, by hand with absent°: the constant
;; may never occur in a variable that existed before it.
(show "with absent°: exists Q forall x. Q = x"
  (°->list #f
    (fresh° (q)
      (let ((c (make-eigen 1)))
        (and° (absent° c q) (=° q c))))))
(show "with absent°: forall x exists Y. Y = x, then Q = ok"
  (°->list #f
    (fresh° (q)
      (let ((c (make-eigen 1)))
        (and° (absent° c q) (fresh° (y) (=° y c)) (=° q 'ok))))))
(show "with absent°: escape through an alias is caught"
  (°->list #f
    (fresh° (q)
      (let ((c (make-eigen 1)))
        (and° (absent° c q)
              (fresh° (y) (=° q (list y)) (=° y c)))))))
\end{scheme}
\begin{answers}
;; naive: exists Q forall x. Q = x
((λ () (make-record-instance 'eigen 1)))
;; with absent°: exists Q forall x. Q = x
()
;; with absent°: forall x exists Y. Y = x, then Q = ok
((λ () 'ok))
;; with absent°: escape through an alias is caught
()
\end{answers}
Three cautions come from the notes on \micro{}. First, \scm{μkanren-state-find/value} rebuilds records,
so eigens must be compared by an identifier (with \scm{equal?}), never by \scm{eq?}; and two eigens
must never share an identifier. Second, \micro{} binds the \emph{older} variable to the \emph{newer}
one in a variable-variable equation. So \enquote{rank = creation index} is not a scope discipline by
itself, and creation time is not scope depth. Third, the check is first-order only: there is no
higher-order unification here.
\end{kanren}

\subsection{Moiseenko 2019: constructive negation}

Moiseenko \cite{Moiseenko2019} builds the answers to $\neg g$ by collecting \emph{all} answers of $g$
and negating them. To make this closed under negation, states become formulas (eq.~(1), p.\,6):
\[
  \exists\bar x.\Big(\bigwedge_i v_i=t_i\ \wedge\ \bigwedge_j\forall\bar y_j.\bigvee_k w_{jk}\neq u_{jk}\Big).
\]
This needs a solver for \emph{universally quantified disequalities}:
\begin{lemma}[{\cite[Lemma 3, p.\,9]{Moiseenko2019}}]
$t\neq u$ is satisfiable iff $\mathrm{mgu}(t,u)\neq\top$.
\end{lemma}
\noindent This is exactly what \micro{}'s \scm{≠°} implements.
\begin{lemma}[{\cite[Def.~11, Lemma 4, p.\,10]{Moiseenko2019}}]
Let $\theta=\mathrm{mgu}(t,u)$, oriented so that $\bar x\not\subseteq\mathrm{codom}\,\theta$. Then
$\forall\bar x.\,t\neq u$ is satisfiable iff $\theta=\bot$ or $\mathrm{dom}(\theta)\not\subseteq\bar x$,
and unsatisfiable iff $\mathrm{dom}(\theta)\subseteq\bar x$.
\end{lemma}
\noindent The orientation is a normalization: if the mgu maps an outer variable $v$ to some $x\in\bar x$,
flip that binding to $x\mapsto v$. In words: $\forall\bar x.\,t\neq u$ is unsatisfiable exactly when $t$
and $u$ can be made equal by choosing values for the universal variables alone. Example: for
$\forall x.\,x\neq\alpha$, the mgu is $\{x\mapsto\alpha\}$, its domain is $\{x\}\subseteq\bar x$, so the
constraint is unsatisfiable. For $\forall x.\,\alpha\neq(1,x)$ the mgu $\{\alpha\mapsto(1,x)\}$ binds the
outer $\alpha$, so it is satisfiable (by any $\alpha$ that is not a pair with first component~$1$).
\begin{lemma}[{\cite[Def.~12, Lemma 5, p.\,10--11]{Moiseenko2019}}]
For $\forall\bar x\exists\bar y.\,t\equiv u$, let $\theta$ be idempotent and $\hat\theta$ its restriction
to the non-$\bar y$ variables. The formula is satisfiable iff $\mathrm{dom}(\hat\theta)\cap\bar x=\emptyset$
and no $p\in\mathrm{codom}(\hat\theta)$ has $\mathrm{Vars}(p)\cap\bar x\neq\emptyset$.
\end{lemma}
\noindent The notes observe that Lemma 5 is exactly the eigenvariable condition. Universal variables
may not be bound, and they may not appear in the bindings of outer variables. With negation one gets
$\mathit{ifte}\;c\;t\;e=(c\wedge t)\vee(\neg c\wedge e)$, $g_1\Rightarrow g_2=\neg g_1\vee(g_1\wedge g_2)$
and $\mathit{forall}\;g=\neg\,\mathit{fresh}\,(x)\,(\neg\,g\,x)$. The price: \emph{stratification} is
required, negating a goal with infinitely many answers diverges, and infinitely many constructors of
each arity are assumed. A remark on p.\,4: \enquote{there is no sound implementation of eigen variables
with the support of disequality constraints. We observed that our implementation of universal
quantifications through double negation works nicely with disequalities}.

\begin{pitfall}[eigen-constants and disequality]
λProlog's generic ∀ is sound because λProlog has no negation. Add \scm{≠°} and consider
$\forall x.\,x\neq\alpha$, with $\alpha$ an outer variable. It must \emph{fail}: take $x:=\alpha$.
Moiseenko's Lemma~4 agrees, since the mgu $\{x\mapsto\alpha\}$ has domain $\subseteq\{x\}$. A naive
eigen-constant $c$ makes $c\neq\alpha$ hold forever, because $\alpha$ can never be bound to $c$, and
the query wrongly succeeds. The same issue arises for type constraints: $c$ stands for an arbitrary
term, so \scm{symbol°} of $c$ should not be provable unless assumed. It also arises for
\scm{absent°}.
\end{pitfall}

\subsection{Jin, Rosenblatt, Might and Zhang 2021: ∀ and ⇒ by relative complement}

Jin et al.\ \cite{JinEtAl2021} build on the \emph{first-order miniKanren representation}
\cite{RosenblattEtAl2019}, where goals are data that can be inspected. \scm{(forall (v) domain goal)}
works in four steps: it checks the domain's satisfiability; it extracts candidates by running
\scm{(exists (v) (conj domain body))} and splitting the states into disjunction-free ones; it
computes a \emph{relative complement}, recursing on \scm{(forall (v) (conj relcomp domain) body)};
and it applies a domain relative negation. This is an \emph{extensional} ∀ over the universe of
literal values: the classical \enquote{for every value} reading, not λProlog's generic one.
\scm{(implies A B)} combines \emph{semantic} solving, which rewrites to $\neg A\vee B$ when $A$ is
negatable, with \emph{syntactic} solving. In the latter, an assumptions list is threaded through, a
user-relation call can be solved by an assumed call, and relations are unfolded under a
\emph{closed-world} reading. The paper's opening examples (p.\,1):

\begin{scheme}
> (run 1 (a) (forall (v) () (== v a)))                                   '()
> (run 1 () (forall (a) () (exists (v) () (== v a))))                    '(())
> (run 1 (a b) (forall (v) () (disj* (== v a) (=/= v b))))              '((_.0 _.0))
> (run 1 (b) (forall (a) () (disj* (not-symbolo a) (=/= a b))))         '((#f))
> (run 1 (f) (implies (evalo f 1) (exists (a) (evalo f a))))            '((_.0))
\end{scheme}

\noindent The third query is the instructive one. $\forall v.\,(v=a\vee v\neq b)$ holds when $a=b$.
That is a classical case split on $v$, which λProlog's generic ∀ cannot do; λProlog has no $\neq$ at
all. The authors list the limitations themselves (§7): the system is slow, produces duplicates, is
incomplete, and soundness is unproven. Their related work notes that λKanren 2020 supports only
relation calls and equalities in antecedents, with \enquote{no unfolding \dots\ consistent with an
open-world interpretation}.

\begin{pitfall}[three different ∀]
λProlog's ∀ is \rulename{generic}: intuitionistic, proved by a new constant, with no case analysis
(\cref{ch:hh}). Jin et al.'s ∀ is extensional: \enquote{for every value in the universe}, by relative
complements. Moiseenko's ∀ is $\neg\exists\neg$: classical, stratified, and it needs finitely many
answers. A query that succeeds in one may fail or diverge in another. For instance, Jin et al.'s §6.6
shows that \scm{(forall (expr) () (disj (neg (halto expr)) (halto expr)))} does not terminate. Only the
first ∀ is the book's.
\end{pitfall}

% =====================================================================================================
\section{How λProlog systems do it}\label{sec:lkanren-impl}

\subsection{Miller 1991: \texorpdfstring{$\Llam$}{L-lambda} by transitions over a mixed prefix}\label{sec:lkanren-miller}

\Cref{ch:hounif} presents pattern unification as the book does. Here we recall only the shape of
Miller's original formulation \cite{Miller1991}, because both λKanren 2020 and Nipkow transcribe it.
The book's statement of the restriction is: \enquote{a program clause or goal formula $B$ is in
$\Llam$ if all quantification in $B$ is over nonpredicate types and every subterm of $B$ of the form
$(x\;t_1\dots t_n)$ in which $x$ is an essentially existentially quantified variable is such that
$t_1,\dots,t_n$ are distinct variables that are universally quantified within the scope of the
quantifier binding $x$} \book{220}. Then only β$_0$ is needed: $(\abs{y}{s})\,x=s[x/y]$ when $x$ is a
variable not free in $\abs{y}{s}$.

Miller's interpreter turns \rulename{generic} into a \emph{meta-level} universal quantifier, and
\rulename{backchain} turns the clause's ∀s into ∃s placed at the point of use. This is why mixed
prefixes $\forall\exists\forall\dots$ arise dynamically. The unification transitions are ξ (descend
under binders, renaming apart), rigid-rigid, \emph{raising}, \emph{pruning}, flex-flex and flex-rigid.
The algorithm terminates (Theorem 7.1), and when it succeeds it yields a most general unifier
(Proposition 7.3). One of Miller's worked examples (p.\,14--15):
\[
  Q\exists u\exists v\forall y\,[f(\abs{x}{f(u\,x\,y)})=f(\abs{w}{v\,y})]
  \xrightarrow{\text{rigid}}\abs{x}{f(u\,x\,y)}=\abs{w}{v\,y}
  \xrightarrow{\xi}\forall y\forall x\,[f(u\,x\,y)=v\,y]
\]
\[
  \xrightarrow{\text{pruning}}[u\mapsto\abs{x}{\abs{y}{u'\,y}}]
  \xrightarrow{\text{flex-rigid}}[v\mapsto\abs{y}{f(u'\,y)}].
\]
Read the last two steps with the prefix in mind, which is now $Q\exists u\exists v\forall y\forall x$. The
equation is flex-rigid: $v\,y$ against $f(u\,x\,y)$. The universal $x$ is quantified after $v$ and is not
among $v$'s arguments, so $v$'s solution may not mention it. Every occurrence of $x$ sits under the
flexible head $u$, so the problem is not hopeless: \emph{pruning} removes $u$'s dependency on its first
argument, $u\mapsto\abs{x}{\abs{y}{u'\,y}}$. After that the rigid side is $f(u'\,y)$, and flex-rigid
abstracts $v$'s argument: $v\mapsto\abs{y}{f(u'\,y)}$. Had $x$ occurred \emph{outside} any flexible
head, pruning would fail and so would the equation.

In the untyped version (§9.3), rigid-rigid with different arities fails, and η-expansion is used only
on terms that must have functional type. This matters for an untyped Scheme term language. Miller also
shows how to stay inside $\Llam$ by replacing meta-level β with \emph{copy clauses} and hypothetical
reasoning, as in the book's Figure~7.10 \book{200}:

\begin{lprolog}
subst M T S :- pi x\ copy x T => copy (M x) S.
\end{lprolog}

\noindent The same trick encodes a \emph{non-pattern} problem inside $\Llam$ \cite[§10.3, p.\,30--31]{Miller1991}.
With copy clauses \lp{copyterm a a}, \lp{copyterm (f X) (f U) :- copyterm X U.} and so on for the
signature, define
\begin{lprolog}
substterm M T N :- pi c\(copyterm c T => copyterm (M c) N).
\end{lprolog}
\noindent The unifiers of $F\,a\lkq g\,a\,a$, which is not a pattern because $a$ is a constant, are exactly
the $F$ that make \lp{substterm F a (g a a)} provable. Miller lists the four answers an $\Llam$
interpreter should return: \lp{w\(g w w)}, \lp{w\(g w a)}, \lp{w\(g a w)}, \lp{w\(g a a)}. The choice
\enquote{is this $a$ the argument or the constant?} is made by backchaining, either on the hypothesis
\lp{copyterm c a} or on the clause \lp{copyterm a a}. Huet's branching has become ordinary proof
search, which a miniKanren would interleave.

\noindent The book's position on programs outside $\Llam$ is pragmatic. A common resolution is
\enquote{to not impose the $\Llam$ restriction on program clauses statically but rather to expect the
condition to be satisfied by the unification problems that are encountered dynamically}. One
\enquote{might defer further consideration of this equation until such time that substitutions
resulting from solving other parts of the problem have altered its status} \book{225--226}.

\subsection{Nipkow 1993: functional pattern unification on de Bruijn terms}

Nipkow \cite{Nipkow1993} restricts attention to $\forall\exists\forall$ problems, where raising is never
needed. Existential variables are global (\texttt{V F}), outer universals are constants (\texttt{C c}),
and inner universals are loose de Bruijn indices (\texttt{B i}). The substitution is triangular. From
the author's ML code, the algorithm is:
\begin{enumerate}
  \item head-normalize both sides through the substitution, with β$_0$ done by renaming indices;
  \item $\lambda s=\lambda t\Rightarrow s=t$; η on demand, $\lambda s=t\Rightarrow s=(\uparrow t)\,0$;
  \item rigid-rigid: equal heads give pairwise unification, different heads fail;
  \item flex-rigid $F\,\bar y=t$: occurs check, then \emph{invert} $t$ with respect to $\bar y$ and
    prune the flexible subterms that mention indices outside $\bar y$;
  \item flex-flex, same head: $F\,\bar y=F\,\bar z\Rightarrow F:=\lambda^{|\bar y|}.H\,(\text{positions where }y_i=z_i)$;
  \item flex-flex, different heads: if $\bar\imath\subseteq\bar\jmath$ then $G:=\lambda\bar\jmath.F\,\bar\imath$,
    symmetrically in the other case, and otherwise both are bound to a fresh $H$ applied to
    $\bar\imath\cap\bar\jmath$.
\end{enumerate}

\subsection{Dowek, Hardin, Kirchner, Pfenning 1996: explicit substitutions}\label{sec:lkanren-dhkp}

DHKP \cite{DowekHardinKirchnerPfenning1996} work in the λσ-calculus, where substitution $a[s]$ is
syntax. Meta-variables are \emph{grafted}, i.e.\ replaced first-order. To recover the logical,
capture-avoiding reading, terms are \textbf{pre-cooked}:
$F(\lambda a,n)=\lambda F(a,n+1)$, $F(a\,b,n)=F(a,n)\,F(b,n)$, $F(k,n)=1[\uparrow^{k-1}]$,
$F(X,n)=X[\uparrow^n]$. A metavariable under $n$ binders is explicitly \emph{shifted past them}. This
is precisely the step that makes grafting agree with substitution, and that static closures omit
(counter-example~1). In a λσ-pattern, every $X[s]\,b_1\dots b_q$ has $s=a_1\dots a_p\cdot\uparrow^n$
with the $a$'s and $b$'s distinct indices. Pattern substitutions are injective and have a right
inverse, so $X[s]=b$ is solved by $X:=b[s^{-1}]$ when every index of $b$ occurs in $s$ (Prop.~8,
inversion). Unification of λσ-patterns is decidable and unitary (Theorem~17), and \enquote{our
presentation does not need mixed prefixes}: raising is just a rule. Outside the fragment, §5.4 says
that equations \enquote{should be postponed as constraints in the hope that they will be instantiated
further}, which \enquote{avoids all branching during constraint simplification}. Empirically,
\enquote{almost all dynamically arising equations fall within the pattern fragment} (p.\,2).

\subsection{Universe indices: Nadathur, Jayaraman and Kwon 1995}

The idea behind every efficient implementation of \rulename{generic} is numeric
\cite{NadathurJayaramanKwon1995}: \enquote{Each time a universal quantifier is encountered, a new
constant must be introduced, giving rise to the next universe in the hierarchy. This requirement can be
accounted for by increasing the “universe index” by 1 and introducing a new constant tagged with this
index. \dots\ When an existential quantifier is encountered, it is instantiated by a logic variable.
This variable may be tagged with the current value of the universe index, the intended interpretation
of the tag being that a term may be substituted for the variable only if all the constants appearing in
the term have a smaller or identical tag} (p.\,15--16).

\begin{definition}[Binding with universe indices {\cite[p.\,16]{NadathurJayaramanKwon1995}}]\label{def:lkanren-ui}
To bind a variable $X$ with tag $i$ to $t$:
(1)~perform the usual occurs check; (2)~check that $t$ contains no constant with tag $>i$; (3)~\emph{before binding}, change to $i$ the tags of the variables of $t$ whose tag is $>i$.
\end{definition}

Step (3) is \emph{lowering}. It is \enquote{required in order to prevent a later instantiation of
these variables from violating the restrictions on instantiations for X}.

\begin{tutorial}[why lowering is needed]
Program: $\forall z(q(z)\imp p(d(z)))$. Goal: $\exists x\forall y(q(y)\imp p(x))$.
\begin{enumerate}
  \item $\exists x$ creates $X$ with tag~1. $\forall y$ raises the index to 2 and creates $c$ with
    tag~2. \rulename{augment} adds $q(c)$. The goal is now $p(X^1)$.
  \item \rulename{backchain} on the clause gives a new variable $Z$ tagged with the current index, 2.
    Unifying $p(X^1)=p(d(Z^2))$ binds $X:=d(Z)$, and \textbf{lowers} $Z$ to tag~1.
  \item The subgoal $q(Z^1)$ meets the hypothesis $q(c^2)$. Binding $Z^1:=c^2$ violates step~(2), so
    it fails.
\end{enumerate}
Without lowering, $Z^2:=c$ would be accepted and $X$ would become $d(c)$. That smuggles the
eigen-constant into a variable quantified outside its scope. The failure is correct: the goal would
need $x=d(y)$, but $x$ is quantified before $y$.
\end{tutorial}

The abstract interpreter works on tuples $\langle G,P,I\rangle$. A ∃-goal tags its new variable with
$I$; a ∀-goal goes to $\langle[c/x]G,P,I{+}1\rangle$; and backchaining instantiates the clause's
variables with tag $I$. The interpreter is sound and complete with respect to intuitionistic
provability, and goal selection order is irrelevant (Prop.~4.1). In the WAM, a \emph{UI register} is
incremented on entering a universal goal and decremented on its completion, and choice points save
it.

\subsection{Teyjus: suspensions and on-the-fly raising}

Teyjus represents term-bound variables by de Bruijn indices and logic variables by mutable cells
\cite{Nadathur2005}. β-reduction is made lazy with the \emph{suspension notation}
$[[t,ol,nl,e]]$: \enquote{a term $t$ that used to occur inside $ol$ abstractions but now appears within
$nl$ of them}, with an environment $e$. Some of the rules:
\begin{gather*}
  (\beta_s)\ (\lambda t_1)\,t_2\to[[t_1,1,0,(t_2,0)::nil]]\qquad
  (r_9)\ [[t,0,0,nil]]\to t\\
  (r_7)\ [[\lambda t,ol,nl,e]]\to\lambda[[t,ol{+}1,nl{+}1,@nl::e]]
\end{gather*}
Terms also carry a closed/open annotation bit; empirically, substitution walks were reduced to
\enquote{a third to an eighth}.

Nadathur and Linnell \cite{NadathurLinnell2005} eliminate the explicit prefix. Each variable gets a
tag $l(y)$, the number of ∃-followed-by-∀ alternations up to its quantifier, and \enquote{a universal
variable $u$ is quantified within the scope of the quantifier for an existential variable $x$ just in
case $l(x)<l(u)$}. Raising is restricted to the universals that actually appear as arguments. A single
traversal of the rigid side (\texttt{mksubst}/\texttt{bnd}) computes the binding, i.e.\ the inversion,
and does pruning and raising on the fly. A non-pattern subterm met during the traversal can be replaced
by a fresh $h(\bar y)$ with the equation $h(\bar y)=t$ \textbf{deferred}. The book describes the same
idea of \enquote{attaching numerical indices to the quantified variables} \book{225}. It also notes that
Teyjus~2 implements only $\Llam$ unification \book{228}.

\begin{example}[no mgu outside the fragment {\cite[p.\,5]{NadathurLinnell2005}}]\leavevmode\\
The problem $\forall u_1\forall u_2\exists x\,(x(u_2)=u_1(u_2))$ is not a pattern: $u_2$ is quantified
\emph{outside} the scope of $x$, so it is a constant for $x$. It has \enquote{two incomparable solutions},
$x\mapsto\abs{w}{u_1(w)}$ and $x\mapsto\abs{w}{u_1(u_2)}$ (written $\lambda(1,u_1(1))$ and
$\lambda(1,u_1(u_2))$ in the paper's de Bruijn notation), and no most general one. Move the quantifier,
$\forall u_1\exists x\forall u_2\,(x(u_2)=u_1(u_2))$, and it becomes a pattern with the single solution
$x\mapsto\abs{w}{u_1(w)}$. A unifier that must return one mgu can only defer the first problem. An
enumerating one would return both answers, as Miller's \texttt{substterm} does for
$F\,a\lkq g\,a\,a$ (\cref{sec:lkanren-miller}).
\end{example}

\subsection{ELPI: de Bruijn levels and a reduction-free fragment}

ELPI \cite{DunchevEtAl2015} starts from an observation about the two roles of β:

\begin{lprolog}
of (lam F) (arr A B) :- pi x\ of x A => of (F x) B.
cbn (app (lam F) N) M :- cbn (F N) M.
\end{lprolog}

\noindent In \lp{of}, β only replaces a bound name with a fresh constant. In \lp{cbn}, β is genuine
object-level substitution. The design has three parts.

\paragraph{Levels for the prefix.} \enquote{When a clause's head is unified with the goal in a context
of length $n$, the universally quantified variables of the clause are instantiated to unification
variables $X^n$ where the level $n$ records that $X$ has only visibility of the initial prefix of length
$n$ of the context. If later a fresh constant is added by the pi rule, the constant occupies position
$n+1$ \dots\ and it will not be allowed to occur in instances of the variable $X^n$} (p.\,4). Level
lowering is done by pruning: the higher-level variable is bound to a fresh lower-level one.

\paragraph{De Bruijn levels, not indices.} The example is $\abs{x}{(\abs{y}{\abs{z}{f\,x\,y\,z}})\,x}$:
\begin{align*}
  \text{indices:}\quad& \lambda x.(\lambda y.\lambda z.f\,x_2\,y_1\,z_0)\,x_0\to_\beta\lambda x.\lambda z.f\,x_1\,x_1\,z_0
    && (x \text{ is renumbered}),\\
  \text{levels:}\quad& \lambda x.(\lambda y.\lambda z.f\,x_0\,y_1\,z_2)\,x_0\to_\beta\lambda x.\lambda z.f\,x_0\,x_0\,z_1
    && (x \text{ keeps its number}).
\end{align*} Three properties follow. A variable bound in Γ keeps its name in any extension
of Γ (DBL1) and in a β-reduct (DBL2), and pushing a binder into the context needs no lifting (DBL3).
\enquote{variables already pushed in the context are treated exactly as constants, and \dots\ the two
notions of level --- De Bruijn's and the position in the context --- coincide}. Terms are
$t::=x_i\mid X^j\mid\lambda x_i.t\mid t\,t$, with global constants as negative $i$; $X^j$ sees the names
below $j$.

\begin{definition}[Reduction-free fragment $L^\beta_\lambda$ {\cite[Def.~1]{DunchevEtAl2015}}]
Every occurrence of a unification variable $X^j$ is applied to $x_j\dots x_{j+k-1}$, for $k\ge0$: exactly
the next $k$ context names, in order.
\end{definition}

\noindent This fragment is contained in $\Llam$ (Property 1). Head reductions of
$(X^j\,x_j\dots x_{j+k-1})\sigma$ take constant time (Property 2), and so does solving
$X^j\,x_j\dots x_{j+k-1}\equiv t$ when no occurs check is needed and the maximum level of $t$ is cached
(Property 3). But $L^\beta_\lambda$ is \textbf{not stable under resolution} (p.\,7), so the fast path
must fall back to general pattern unification. ELPI reports large speed-ups: typeof in 0.27\,s versus
5.64\,s for Teyjus, reduce-cbv in 0.15\,s versus 11.11\,s.

\paragraph{Outside the fragment,} ELPI's default is an \textbf{error}: \enquote{Unification problem
outside the pattern fragment. \dots\ Pass -delay-problems-outside-pattern-fragment \dots\ in order to
delay (deprecated, for Teyjus compatibility)}. Delay is offered instead through \emph{syntactic
constraints}. \lp{declare_constraint (even X) [X]} suspends a goal on variables, and the goal is resumed
when any of them is assigned, so \lp{declare_constraint (even X) [X], X = 1.} fails. Suspended goals keep
their hypothetical context. The documentation's example (\texttt{ELPI.md}, \enquote{Syntactic
constraints}):
\begin{answers}
goal> pi x\ sigma Y\ even x => declare_constraint (even Y) [Y].
Success:
Constraints:
 {x} : even x ?- even (W x)  /* suspended on W */
\end{answers}
\noindent The residual constraint remembers the eigenvariable \texttt{x} and the hypothesis
\lp{even x}. Note also the automatic raising: \texttt{Y} was created inside the scope of \texttt{x}, so
it is printed as \texttt{W x}. Constraint Handling Rules can operate on suspended goals.

\subsection{Makam}

Makam \cite{StampoulisChlipala2018} is a λProlog dialect with polymorphism, used to prototype type
systems. Two idioms are relevant for us. \texttt{bindmany} is a binder over a \emph{list} of
variables, with \texttt{openmany}/\texttt{applymany} to open many binders at once. And reflective
predicates such as \texttt{refl.isunif} test whether a term is an uninstantiated unification
variable. That test is non-logical, and is the analogue of \micro{}'s \scm{project°}.

\begin{kanren}[what the implementations ask of a miniKanren]
The real systems agree on three mechanisms:
(i)~a \emph{level} on every logic variable and eigen-constant, with the binding check and lowering
    of \cref{def:lkanren-ui}; (ii)~pattern unification with inversion and pruning (Nipkow, Nadathur--Linnell, ELPI); (iii)~a store of \emph{delayed} equations or goals, woken on assignment (DHKP §5.4, Teyjus, ELPI's
    syntactic constraints).
\micro{} has no notion of level. Its creation index is a time stamp, not a scope depth: a variable
created after a ∀-scope has closed has a larger index, but it must not see that scope's constant. On
the other hand, the \scm{assocs}-driven re-verification of D is exactly the wake-up hook that a delay
store needs. Search is unchanged: \scm{or°}'s fair interleaving can host Huet's branching. That fits
Miller's remark that a breadth-first search could compute the converse of his \texttt{prenex}
relation, which the depth-first interpretation cannot \cite{Miller1991}. The design itself is in \cref{ch:design}.
\end{kanren}

% =====================================================================================================
\section{Comparison of the design options}\label{sec:lkanren-compare}

\Cref{tab:lkanren-options} sets the four candidate meanings of \scm{=°} on λ-terms side by side. The
rows summarize the analysis of \cref{sec:lkanren-mf21,sec:lkanren-assess,sec:lkanren-nominal,sec:lkanren-impl}. The table records trade-offs and does not rank the options: the choice,
and the reasons for it, are the subject of \cref{ch:design}.

\begin{table}[ht]
\centering\footnotesize
\caption{Design options for higher-order unification in a miniKanren. No option is
chosen here; see \cref{ch:design}.}
\label{tab:lkanren-options}
\renewcommand{\arraystretch}{1.25}
\begin{tabularx}{\textwidth}{>{\raggedright\arraybackslash}p{2.0cm}*{4}{>{\raggedright\arraybackslash}X}}
\toprule
& \textbf{A.\ Pattern unification, de Bruijn + levels} (Miller, Nipkow, Nadathur--Linnell, ELPI)
& \textbf{B.\ Nominal} (Urban--Pitts--Gabbay, αKanren)
& \textbf{C.\ Static closures} (Ma--Friedman 2021)
& \textbf{D.\ Huet pre-unification, enumerated by the stream}\\
\midrule
Fidelity to the book & exact: $\Llam$ is what Teyjus~2 and ELPI run & a different logic (nominal); no
function variables, no β & non-standard grafting semantics: $\abs{a}{X}\lkq\abs{b}{b}$ succeeds, unlike
λProlog & exact for full higher-order unification (pre-unifiers)\\
Decidable / unitary & yes / yes (mgu) & yes / yes, with a freshness context & claimed; no full
algorithm published (a \enquote{preview}) & semi-decidable; possibly infinitely many pre-unifiers\\
New term forms & λ node (indices), eigen-constant with level, flexible application $X\,\bar y$ & nom,
\texttt{tie}, suspension $\pi\cdot X$ & named and nameless binders, levels, closures & as A, plus
imitation and projection\\
New state & a level per variable, current depth, delayed-equation store & freshness store ∇ of $a\#X$ &
a \enquote{context} of closure equations & as A; branching via \texttt{or°}\\
Cost of α & free (syntactic on indices) & swappings on the fly, suspensions & levels & free\\
β & only β$_0$; general β via copy/subst relations, or delayed & none (write \texttt{substo}) &
\enquote{will be powered} (not given) & full\\
∀-goals (\rulename{generic}) & eigen-constant + level check + lowering & a separate mechanism is still
needed: noms are distinct constants, not eigenvariables & \texttt{all} from λKanren 2020 (prefix) & as
A\\
With \texttt{≠°}, \texttt{absent°}, types & eigens need a ∀-aware check (Moiseenko) & αKanren has no
disequality at all & not discussed & as A\\
Effort in \aux{} & medium--high: a Nipkow-size core, plus levels and a delay store & medium: Byrd's
ch.~11 is about 300 lines; a triangular version exists & unclear & high, plus control of infinite
branching\\
\bottomrule
\end{tabularx}
\end{table}

Some combinations are possible and appear in the literature. Option D can sit on top of A as a
separate, explicitly nondeterministic relation; the book itself presents Huet's procedure as a
search over a matching tree (\cref{ch:hounif}), and miniKanren's interleaving could explore it
fairly. Option B can live beside A, for users who want names as data. Option C behaves differently
from all the others on the three counter-examples of \cref{sec:lkanren-counter}; any implementation
should state its answers to them explicitly. The smaller alternatives for ∀ alone are the explicit
prefix of λKanren 2020 (simplest, but each check walks the prefix and the prefix only grows), the
eigen-absento of Hemann--Byrd--Friedman (smallest change, $O(|B|)$ per check, first-order only), and
the end-of-scope check (weakest). The constructive ∀/⇒/¬ of Moiseenko and Jin et al.\ is orthogonal:
it has a classical, extensional reading and does not implement the book's \rulename{generic} and
\rulename{augment}.

% =====================================================================================================
\clearpage
\section{Summary}\label{sec:lkanren-summary}

\begin{itemize}
  \item Ma and Friedman 2021 \cite{MaFriedman2021} is a workshop \emph{preview}. It has two grammars,
    the operations $\mathrm{ext}$ and $\mathrm{db}$, three examples, and no algorithm, rules or
    theorems.
  \item Its first complaint is that pattern unification rejects $\abs{a}{X}\lkq\abs{b}{b}$. That is
    the scope restriction, which λProlog needs for the soundness of ∀. Its fix is to graft de Bruijn
    levels, giving $X\mapsto\texttt{d0}$. λProlog's own answer is to make the dependency explicit:
    $\abs{a}{(X\,a)}\lkq\abs{b}{b}$ is a pattern with mgu $X\mapsto\abs{x}{x}$ (Miller's
    \enquote{discharging a constant}).
  \item Its second complaint is that binder-merging plus naive reification makes
    $\abs{a}{\abs{a}{X}}\lkq\abs{c}{\abs{d}{c}}$ succeed, so \enquote{true = false}. Static closures
    $\lkclo{\phi}{t}$ make shadowed levels unavailable, and the equation fails. Flex-flex closure
    equations are postponed in a \emph{context}.
  \item The assessment uses three counter-examples. The conjunction separates static closures from
    nominal unification. The dangling \texttt{d0} at top level shows that $X$ has no
    depth-independent meaning. \enquote{True = false} also fails under pattern and nominal
    unification, so it is a bug of binder-merging only. Static closures are grafting on de Bruijn
    terms without DHKP's pre-cooking, and α is no longer a congruence on open terms.
  \item λKanren 2020 transcribes Miller 1991 with an explicit prefix and gensym renaming. Non-patterns
    fail. Hypotheses must be threaded in the state, because a host parameter loses them under delayed
    relation bodies.
  \item Nominal unification (αKanren) handles α with swappings and freshness, using capturing
    substitution. It has no function variables, no β and no disequality, and noms are not
    eigenvariables.
  \item Eigenvariables in miniKanren: a fresh constant without a scope check is unsound. Hemann, Byrd
    and Friedman implement the check as an absento over the earlier variables. Disequality and
    eigen-constants conflict (Moiseenko). Jin et al.'s ∀ is extensional, not generic.
  \item Real λProlog systems use universe indices or levels with lowering (Nadathur et al., Teyjus,
    ELPI), pattern unification with inversion and pruning (Nipkow, Nadathur--Linnell), and delayed
    non-pattern equations (DHKP, Teyjus, ELPI's syntactic constraints). Teyjus adds the suspension
    calculus; ELPI uses de Bruijn levels and the $L^\beta_\lambda$ fast path.
  \item The options are pattern+levels, nominal, static closures, and Huet enumerated by the stream
    (\cref{tab:lkanren-options}). The choice is made in \cref{ch:design}.
\end{itemize}
